
\documentclass[11pt,oneside]{article}

\usepackage[T1]{fontenc}
\usepackage[utf8]{inputenc}
\usepackage{lmodern}
\usepackage{microtype}
\usepackage{geometry}
\geometry{margin=0.92in,headheight=14pt}
\usepackage{amsmath,amssymb,amsthm,mathtools}
\usepackage{booktabs,longtable,tabularx,array}
\usepackage{enumitem}
\usepackage{xcolor}
\usepackage{xurl}
\usepackage{hyperref}
\usepackage{listings}
\usepackage{fancyhdr}
\usepackage{titlesec}
\usepackage{needspace}

\definecolor{proofblue}{RGB}{25,82,130}
\definecolor{newgreen}{RGB}{20,105,70}
\definecolor{openorange}{RGB}{166,87,0}
\definecolor{compviolet}{RGB}{103,55,145}
\definecolor{quietgray}{RGB}{82,90,99}
\definecolor{codegray}{RGB}{248,248,248}
\definecolor{rulegray}{RGB}{95,102,110}

\hypersetup{
  colorlinks=true,
  linkcolor=proofblue,
  citecolor=newgreen,
  urlcolor=proofblue,
  pdftitle={Bhargava Factorials for the 2--3 Semigroup and the Alaoglu--Erdos Conjecture},
  pdfauthor={OpenAI GPT-5.6 Pro},
  pdfsubject={A nonduplicative progress report on p-orderings, generalized factorials, and integer-valued interpolation}
}
\urlstyle{same}

\pagestyle{fancy}
\fancyhf{}
\lhead{Alaoglu--Erd\H{o}s: Bhargava-factorial progress}
\rhead{August 23, 2026}
\cfoot{\thepage}
\renewcommand{\headrulewidth}{0.4pt}

\titleformat{\section}{\Large\bfseries}{\thesection}{0.75em}{}
\titleformat{\subsection}{\large\bfseries}{\thesubsection}{0.75em}{}
\titleformat{\subsubsection}{\normalsize\bfseries}{\thesubsubsection}{0.75em}{}

\setlength{\parindent}{0pt}
\setlength{\parskip}{0.52em}
\setlength{\emergencystretch}{3em}
\setlist[itemize]{leftmargin=2em,itemsep=0.25em,topsep=0.35em}
\setlist[enumerate]{leftmargin=2.2em,itemsep=0.25em,topsep=0.35em}
\setlist[description]{leftmargin=3.4cm,labelsep=0.6cm,itemsep=0.3em,topsep=0.35em}

\newtheorem{theorem}{Theorem}[section]
\newtheorem{proposition}[theorem]{Proposition}
\newtheorem{lemma}[theorem]{Lemma}
\newtheorem{corollary}[theorem]{Corollary}
\newtheorem{conjecture}[theorem]{Conjecture}
\newtheorem{criterion}[theorem]{Proof criterion}
\newtheorem{problem}[theorem]{Research problem}
\theoremstyle{definition}
\newtheorem{definition}[theorem]{Definition}
\newtheorem{experiment}[theorem]{Exact computation}
\newtheorem{warning}[theorem]{Caution}
\theoremstyle{remark}
\newtheorem{remark}[theorem]{Remark}

\newcommand{\N}{\mathbb N}
\newcommand{\Nzero}{\mathbb N_0}
\newcommand{\Z}{\mathbb Z}
\newcommand{\Q}{\mathbb Q}
\newcommand{\F}{\mathbb F}
\newcommand{\R}{\mathbb R}
\newcommand{\Qp}{\mathbb Q_p}
\newcommand{\Zp}{\mathbb Z_p}
\newcommand{\B}{\mathcal B}
\newcommand{\Cout}{\mathcal C}
\newcommand{\GammaTwoThree}{\Gamma_{2,3}}
\newcommand{\Int}{\operatorname{Int}}
\newcommand{\ord}{\operatorname{ord}}
\newcommand{\lcm}{\operatorname{lcm}}
\newcommand{\vp}{v_p}
\newcommand{\cl}{\operatorname{cl}}
\newcommand{\dist}{\operatorname{dist}}
\newcommand{\Span}{\operatorname{span}}
\newcommand{\supp}{\operatorname{supp}}
\newcommand{\rad}{\operatorname{rad}}
\newcommand{\vol}{\operatorname{vol}}
\newcommand{\sgn}{\operatorname{sgn}}
\newcommand{\floor}[1]{\left\lfloor #1\right\rfloor}
\newcommand{\ceil}[1]{\left\lceil #1\right\rceil}
\newcommand{\statusproved}{\textcolor{newgreen}{\textbf{[proved here]}}}
\newcommand{\statusexternal}{\textcolor{quietgray}{\textbf{[external theorem]}}}
\newcommand{\statuscomputed}{\textcolor{compviolet}{\textbf{[exact computation]}}}
\newcommand{\statusopen}{\textcolor{openorange}{\textbf{[open target]}}}
\newcommand{\statusconditional}{\textcolor{openorange}{\textbf{[conditional criterion]}}}

\newenvironment{statusbox}
{\begin{center}\begin{minipage}{0.94\textwidth}\raggedright\hrule
\vspace{0.42em}\textbf{Status.}\ }
{\par\vspace{0.32em}\hrule\end{minipage}\end{center}}

\newenvironment{resultbox}
{\begin{center}\begin{minipage}{0.94\textwidth}\hrule\vspace{0.45em}}
{\vspace{0.35em}\hrule\end{minipage}\end{center}}

\lstdefinestyle{pythonstyle}{
  language=Python,
  basicstyle=\ttfamily\scriptsize,
  backgroundcolor=\color{codegray},
  frame=single,
  rulecolor=\color{rulegray},
  numbers=left,
  numberstyle=\tiny\color{rulegray},
  numbersep=7pt,
  showstringspaces=false,
  breaklines=true,
  columns=fullflexible,
  keepspaces=true,
  tabsize=4
}

\lstdefinestyle{cppstyle}{
  language=C++,
  basicstyle=\ttfamily\scriptsize,
  backgroundcolor=\color{codegray},
  frame=single,
  breaklines=true,
  columns=fullflexible,
  keepspaces=true,
  showstringspaces=false,
  numbers=left,
  numberstyle=\tiny\color{quietgray},
  numbersep=6pt
}


\lstdefinestyle{leanstyle}{
  basicstyle=\ttfamily\scriptsize,
  backgroundcolor=\color{codegray},
  frame=single,
  breaklines=true,
  columns=fullflexible,
  keepspaces=true,
  showstringspaces=false,
  tabsize=2
}

\begin{document}
\pagenumbering{gobble}

\begin{titlepage}
\centering
\vspace*{0.7cm}
{\Huge\bfseries Bhargava Factorials for the \(2\)--\(3\) Semigroup\par}
\vspace{0.35cm}
{\LARGE and a New Integer-Valued-Interpolation Frontier\par}
\vspace{0.7cm}
{\Large A nonduplicative progress report on the\par}
\vspace{0.12cm}
{\Large Alaoglu--Erd\H{o}s integer-exponent conjecture\par}
\vspace{1.0cm}
{\large \(2^x,3^x\in\Z\Longrightarrow x\in\Z\)\par}
\vfill

\begin{minipage}{0.92\textwidth}
\small
\textbf{Bottom line.}
No complete proof is claimed.  The attached unified report already develops an exceptionally
large collection of transcendence, determinant, interpolation, \(p\)-adic, dynamical, and
formalization arguments.  This report deliberately enters the one substantial interface that
the supplied corpus names but does not develop: Bhargava \(p\)-orderings and generalized
factorials.

For the rank-two semigroup
\[
  \B=\{2^a3^b:a,b\in\Nzero\},
\]
we determine the complete local factorial calculus.  At every prime \(p\ge5\), if
\[
  h_p(e)=\bigl|\langle2,3\rangle\subset(\Z/p^e\Z)^\times\bigr|,
\]
then
\[
  v_p(k!_{\B})=\sum_{e\ge1}\floor{\frac{k}{h_p(e)}}.
\]
At \(p=3\), the semigroup is dense in \(\Z_3\), so \(v_3(k!_{\B})=v_3(k!)\).  At \(p=2\),
the closure splits self-similarly into an odd orbit and a doubled copy of itself; this gives
an exact merge recursion and the asymptotic slope
\[
 \frac{v_2(k!_{\B})}{k}\longrightarrow\frac{\sqrt{11}-1}{2}.
\]
A collision lemma proves that every prime power contributing to \(k!_{\B}\) is below
\(6^{\floor{\sqrt{k}}}\), yielding a finite exact algorithm.  We also construct the first
regular-basis element that cannot arise from one global Newton ordering and determine the
optimal normalized Vandermonde defects through eight nodes:
\[
  1,1,1,1,1,7,7,396.
\]

These local and finite results feed an optimal normalized interpolation determinant.  Under a hypothetical
nonintegral solution \(x\), the map \(q\mapsto q^x\) is integer-valued on \(\B\); therefore every
regular Bhargava basis produces nonzero integer determinants.  This turns the conjecture into a
concrete subunit problem in the exact lattice \(\Int(\B,\Z)\), and also into a finite-support
adelic problem after localizing at \(2\) and \(3\).  The new machinery removes all ambiguity
about automatic denominator content.  What remains is genuinely Archimedean--nonarchimedean
coupling: one must control cancellation or the denominator carried by the structural primes.
\end{minipage}

\vfill
{\large August 23, 2026\par}
\end{titlepage}

\pagenumbering{roman}

\begin{abstract}
The Alaoglu--Erd\H{o}s conjecture is the assertion that a real number \(x\) for which
\(2^x\) and \(3^x\) are integers must itself be an integer.  If a nonintegral solution existed
and \(M=2^x\), \(N=3^x\), then the power map would define an order-preserving monoid
isomorphism
\[
  2^a3^b\longmapsto M^aN^b
\]
between two rank-two integer semigroups.  The supplied unified report studies this situation
through many complementary frameworks, but leaves Bhargava's theory of \(p\)-orderings as an
undeveloped research interface.

This report computes the Bhargava factorial of the input semigroup \(\B=\langle2,3\rangle_+\).
For odd unramified primes the answer is an exact orbit-size formula; the prime \(3\) is
ordinary because \(\B\) is \(3\)-adically dense; and the prime \(2\) is governed by a
self-referential two-branch renewal law whose limiting slope is
\((\sqrt{11}-1)/2\).  A finite-support theorem bounds every contributing prime power and makes
the factorial exactly computable.  We deduce fixed-prime excess formulas, denominator-growth
bounds, and a rank-two Artin--Bhargava asymptotic problem.

The factorial calculus is then inserted into integer-valued interpolation.  A regular basis of
\(\Int(\B,\Z)\) gives an optimally normalized generalized Vandermonde determinant which, under
a counterexample, is a nonzero integer.  The first non-Newton basis element is given explicitly,
and exhaustive exact certificates yield the defect minima \(1,1,1,1,1,7,7,396\) through
eight nodes.  Equivalently, no degree-\(d\) integer-valued
polynomial can approximate \(t^x\) within one unit at \(d+2\) semigroup nodes.  We formulate
four finite proof criteria, an \(S\)-integer compact-group version, and a Lean-oriented
formalization plan.  The conjecture remains open: generalized factorials settle the exact local
divisibility, but not the output-sensitive cancellation needed to defeat the structural-prime
denominator tax.
\end{abstract}

\tableofcontents
\newpage
\pagenumbering{arabic}

\section{Scope, trust boundary, and relation to the supplied report}

\subsection{The mathematical starting point}

The conjecture, introduced in the study of highly composite and related numbers~\cite{AE1944}, is
\begin{conjecture}[Alaoglu--Erd\H{o}s]
\label{conj:AE}
For \(x\in\R\),
\[
  2^x\in\Z,\qquad 3^x\in\Z
  \quad\Longrightarrow\quad
  x\in\Z.
\]
\end{conjecture}
A negative exponent cannot occur, because then \(0<2^x<1\).  A rational solution is an
integer by unique factorization.  An algebraic irrational solution is excluded by the
Gelfond--Schneider theorem~\cite{Baker1975}.  Thus any counterexample is positive and transcendental.  With
\[
  M=2^x,\qquad N=3^x,
\]
one obtains
\begin{equation}
 (2^a3^b)^x=M^aN^b\in\Z
 \qquad(a,b\in\Nzero).
 \label{eq:semigroup-power}
\end{equation}
The problem is therefore not a two-value accident: a counterexample would make a real power
integer-valued on an entire rank-two multiplicative semigroup.

\subsection{What is intentionally not repeated}

The supplied report is more than twenty-three thousand lines long.  It already treats, among
many other topics, the four-exponentials barrier, exact conditional semigroup structure,
valuation lattices, finite differences, generalized Vandermonde determinants, tropical
divisibility ceilings, Sturmian codings, rational-point counting, Loewner matrices, geometric
Newton interpolation, \(q\)-special functions, Kummer defects, and formalization boundaries.
Repeating those arguments would add bulk rather than information.

The present report follows one explicit opening in that corpus.  Its arithmetic-interpolation
chapter states that broader Bhargava-ordering proposals remain a research interface, but it
does not define a \(p\)-ordering, compute a generalized factorial, or build the corresponding
regular basis.  Everything below is organized around that missing local calculus.

\subsection{Evidence labels}

The labels used here have literal meanings.
\begin{description}
\item[\statusproved] A complete paper proof is included.
\item[\statuscomputed] The assertion is the output of exact integer arithmetic; source is
included and the mathematical cutoff is proved.
\item[\statusexternal] The result is quoted from the literature and is not reproved in full.
\item[\statusconditional] The statement is a rigorously proved sufficient condition, but its
hypothesis has not been established.
\item[\statusopen] The statement is proposed as a research target or conjectural asymptotic.
\end{description}
No numerical pattern is promoted to a theorem.  In particular, the public announcement of a
rival Alaoglu--Erd\H{o}s breakthrough is not available in a citable proof or preprint at the
time of writing and is not used.

\section{Bhargava factorials and the exact interpolation lattice}
\label{sec:bhargava-basics}

\subsection{\(p\)-orderings}

Let \(S\subset\Z\) be infinite and let \(p\) be a prime.  A \(p\)-ordering of \(S\) is a
sequence \(a_0,a_1,\ldots\) chosen recursively so that, after \(a_0,\ldots,a_{k-1}\) have been
chosen, the value
\[
  v_p\!\left(\prod_{i=0}^{k-1}(a-a_i)\right)
\]
is minimized over \(a\in S\).  Bhargava proved that the minimum is independent of the choices
made in the ordering~\cite{Bhargava1997,Bhargava1998}.  Denote it by
\[
  \nu_{S,p}(k).
\]
The \(k\)-th generalized factorial is
\begin{equation}
  k!_S=\prod_p p^{\nu_{S,p}(k)}.
  \label{eq:generalized-factorial}
\end{equation}
For \(S=\Z\), the ordinary ordering \(0,1,2,\ldots\) gives
\[
  \nu_{\Z,p}(k)=v_p(k!),\qquad k!_{\Z}=k!.
\]
If \(S\subset T\), minimization over the smaller set can only increase the valuation, so
\(k!_T\mid k!_S\).

The factorial is not merely a divisibility statistic.  Define
\[
  \Int(S,\Z)=\{P\in\Q[T]:P(S)\subset\Z\}.
\]
Bhargava's factorial and characteristic-ideal theory~\cite{Bhargava1997,Bhargava2000,Bhargava2009,CahenChabert}
implies that, because \(\Z\) is principal, there is a regular basis
\[
  F_0,F_1,F_2,\ldots\in\Int(S,\Z),
  \qquad
  \deg F_k=k,
  \qquad
  \operatorname{lc}(F_k)=\frac{1}{k!_S},
  \label{eq:regular-basis}
\]
and
\[
  \Int_{\le d}(S,\Z)
  =\bigoplus_{k=0}^{d}\Z F_k.
  \label{eq:int-lattice}
\]
Thus \(k!_S\) is the optimal leading-coefficient denominator for degree \(k\).

\subsection{Closure invariance}

The local computation becomes transparent after passing to \(p\)-adic closures.

\begin{lemma}[Closure invariance]
\label{lem:closure-invariance}
Let \(S\subset\Z\) be infinite and let \(\overline S^{\,p}\) be its closure in \(\Z_p\).
Then
\[
  \nu_{S,p}(k)=\nu_{\overline S^{\,p},p}(k)
\]
for every \(k\ge0\), where the right side is defined by the same greedy rule in the compact
set \(\overline S^{\,p}\).
\end{lemma}

\begin{proof}
Fix finitely many previously selected points.  The valuation of their difference product is an
integer-valued locally constant function away from those points.  The minimum over the compact
closure is finite because the closure is infinite.  A minimizer can be approximated by an element
of \(S\) closely enough that every difference has the same valuation.  Hence the minimum over
\(S\) is no larger than the minimum over the closure; the reverse inequality is immediate from
\(S\subset\overline S^{\,p}\).  Induction on the stage gives the claim.
\end{proof}

\subsection{A procyclic orbit lemma}

\begin{lemma}[Orbit factorial]
\label{lem:orbit-factorial}
Let \(H\subset\Z_p^\times\) be a procyclic closed subgroup, let \(g\) be a topological
generator, and let
\[
  h(e)=|\operatorname{im}(H\to(\Z/p^e\Z)^\times)|.
\]
Then \(1,g,g^2,\ldots\) is a \(p\)-ordering of \(H\), and
\begin{equation}
  \nu_{H,p}(k)=\sum_{e\ge1}\floor{\frac{k}{h(e)}}.
  \label{eq:orbit-factorial}
\end{equation}
\end{lemma}

\begin{proof}
For a candidate \(y\in H\),
\[
 v_p\!\left(\prod_{i<k}(y-g^i)\right)
 =\sum_{e\ge1}\#\{0\le i<k:y\equiv g^i\pmod{p^e}\}.
 \label{eq:level-count}
\]
At level \(e\), the first \(k\) powers of \(g\) are distributed among the \(h(e)\) residue
classes with occupancies differing by at most one.  Every candidate therefore meets at least
\(\floor{k/h(e)}\) previous points in its class.  The candidate \(g^k\) attains that minimum
simultaneously at every level, because the preceding indices congruent to \(k\) modulo \(h(e)\)
are \(k-h(e),k-2h(e),\ldots\).  Summing the level occupancies proves the formula.
\end{proof}

\begin{remark}
The formula is local but global in depth.  A single number \(h(e)\) records the entire
two-generator orbit modulo \(p^e\); no enumeration of semigroup elements is required.
\end{remark}

\section{The complete local factorial of \(\B=\langle2,3\rangle_+\)}
\label{sec:local-factorial}

Set
\[
  \B=\{2^a3^b:a,b\in\Nzero\}.
\]
For a prime \(p\ge5\), define
\begin{equation}
  H_{p,e}=\langle2,3\rangle\subset(\Z/p^e\Z)^\times,
  \qquad
  h_p(e)=|H_{p,e}|.
  \label{eq:hp-def}
\end{equation}

\subsection{Every unramified odd prime}

\begin{theorem}[Exact unramified formula]
\label{thm:odd-prime-formula}
For every prime \(p\ge5\) and every \(k\ge0\),
\begin{equation}
  \boxed{\displaystyle
  \nu_{\B,p}(k)=\sum_{e\ge1}\floor{\frac{k}{h_p(e)}}.}
  \label{eq:odd-prime-main}
\end{equation}
Moreover,
\begin{equation}
  h_p(e)=\lcm\bigl(\ord_{p^e}(2),\ord_{p^e}(3)\bigr).
  \label{eq:h-lcm}
\end{equation}
\end{theorem}

\begin{proof}
Modulo \(p^e\), nonnegative powers of \(2\) and \(3\) already contain their inverses, because
the ambient group is finite.  Hence the image of \(\B\) is precisely \(H_{p,e}\).  The closure
of \(\B\) in \(\Z_p^\times\) is the closed subgroup generated by \(2\) and \(3\).
For odd \(p\), \(\Z_p^\times\) is procyclic, and every closed subgroup of a procyclic group is
procyclic.  Lemma~\ref{lem:closure-invariance} and Lemma~\ref{lem:orbit-factorial} give
\eqref{eq:odd-prime-main}.  Since \((\Z/p^e\Z)^\times\) is cyclic, the subgroup generated by
two elements has order the least common multiple of their individual orders, proving
\eqref{eq:h-lcm}.
\end{proof}

\subsection{Prime-power lifting and the local density}

The sequence \(h_p(e)\) has a simple eventual form.  Let \(\lambda_p=h_p(1)\).  There is an
integer \(c_p\ge1\), the saturation depth of the closed subgroup, such that
\begin{equation}
  h_p(e)=\lambda_p\,p^{\max(0,e-c_p)}.
  \label{eq:saturation-depth}
\end{equation}
For the usual non-Wieferich behavior, \(c_p=1\).  The exact formula then becomes
\[
 \nu_{\B,p}(k)
 =c_p\floor{\frac{k}{\lambda_p}}
  +\sum_{j\ge1}\floor{\frac{k}{\lambda_p p^j}}.
\]
In particular the fixed-prime slope exists:
\begin{equation}
 \alpha_p:=\lim_{k\to\infty}\frac{\nu_{\B,p}(k)}{k}
 =\sum_{e\ge1}\frac1{h_p(e)}.
 \label{eq:alpha-p}
\end{equation}

For comparison,
\[
 \lim_{k\to\infty}\frac{v_p(k!)}{k}=\frac1{p-1}.
\]
Because \(h_p(e)\le\varphi(p^e)=(p-1)p^{e-1}\),
\begin{equation}
  \alpha_p\ge\frac{p}{(p-1)^2}
  =\frac1{p-1}+\frac1{(p-1)^2}.
  \label{eq:local-excess-lower}
\end{equation}
Thus every fixed unramified prime contributes strictly more factorial mass on \(\B\) than it
does on all of \(\Z\).

\subsection{The prime \(3\): no gain}

\begin{theorem}[Three-adic density]
\label{thm:p3}
The closure of \(\B\) in \(\Z_3\) is all of \(\Z_3\).  Consequently,
\begin{equation}
  \boxed{\nu_{\B,3}(k)=v_3(k!).}
  \label{eq:p3-main}
\end{equation}
\end{theorem}

\begin{proof}
The integer \(2\) is a primitive root modulo \(3^e\) for every \(e\ge1\).  Every nonzero
residue modulo \(3^e\) has the form \(3^b u\), with \(0\le b<e\) and \(u\) a unit modulo
\(3^{e-b}\).  Choose \(a\) with \(2^a\equiv u\pmod{3^{e-b}}\); then \(2^a3^b\) realizes the
residue.  The zero residue is realized by \(3^e\).  Hence \(\B\) is surjective modulo every
\(3^e\), so its closure is \(\Z_3\).  Closure invariance reduces the factorial to the ordinary
one.
\end{proof}

\subsection{The prime \(2\): a self-similar residue tree}

The prime \(2\) is the genuinely new local feature.  Let
\[
  K=\overline{\B}^{\,2}\subset\Z_2,
  \qquad
  U=\overline{\langle3\rangle}\subset\Z_2^\times.
\]
Every element of \(K\) is either odd or twice another element of \(K\), and the two pieces are
separated modulo \(2\):
\begin{equation}
  K=U\;\sqcup\;2K.
  \label{eq:K-decomposition}
\end{equation}

The orbit sizes of \(U\) are
\[
 |\langle3\rangle\bmod2|=1,\qquad
 |\langle3\rangle\bmod4|=2,\qquad
 |\langle3\rangle\bmod2^e|=2^{e-2}\quad(e\ge3).
\]
Hence Lemma~\ref{lem:orbit-factorial} gives
\begin{equation}
  u_k:=\nu_{U,2}(k)
  =k+\floor{\frac{k}{2}}+v_2(k!).
  \label{eq:u-k}
\end{equation}

\begin{theorem}[Exact two-adic merge recursion]
\label{thm:p2-recursion}
Let
\[
 a_k=\nu_{\B,2}(k).
\]
Set \(a_0=0\).  After selecting the initial odd branch, maintain counts \(r\) and \(s\) of
previously selected points from \(2K\) and \(U\), respectively.  Initially \(r=0,s=1\).
At stage \(k\ge1\), compare
\begin{equation}
  E_r=r+a_r,\qquad O_s=u_s.
  \label{eq:branch-costs}
\end{equation}
Set \(a_k=\min(E_r,O_s)\) and increment the count of a branch attaining the minimum.  Any
tie rule gives the same value sequence.  Equivalently, as multisets,
\begin{equation}
 \{a_k:k\ge0\}
 =
 \{j+a_j:j\ge0\}\;\uplus\;\{u_j:j\ge0\},
 \label{eq:multiset-fixedpoint}
\end{equation}
with both sides listed in nondecreasing order.

The first values are
\begin{equation}
 0,0,1,1,3,4,4,5,7,9,9,10,10,12,13,14,15,18,19,19,20,\ldots.
 \label{eq:a-first}
\end{equation}
\end{theorem}

\begin{proof}
Differences between an odd point and an even point are \(2\)-adic units, so cross-branch
points contribute zero valuation.  Within \(U\), the \(s\)-th local cost is \(u_s\).  Within
\(2K\), every difference acquires one factor of \(2\), and after dividing by \(2\) the local
problem is again \(K\); with \(r\) previous even points the cost is therefore \(r+a_r\).
The global greedy ordering is exactly the merge of these two local cost streams.  Independence
of the resulting \(p\)-sequence from tie choices is Bhargava's invariance theorem.
\end{proof}

\subsection{The renewal slope}

\begin{lemma}[Self-similar merge slope]
\label{lem:merge-slope}
Let \(b_k\) be nondecreasing with \(b_k/k\to\beta>0\).  Suppose a nondecreasing unbounded
sequence \(a_k\) is the ordered multiset merge of
\[
  \{k+a_k:k\ge0\}
  \quad\hbox{and}\quad
  \{b_k:k\ge0\}.
\]
Then \(a_k/k\) converges to the unique positive \(c\) satisfying
\begin{equation}
  \frac1c=\frac1{1+c}+\frac1\beta.
  \label{eq:merge-slope-equation}
\end{equation}
\end{lemma}

\begin{proof}
Let \(F(t)=\#\{k:a_k\le t\}\).  The merge identity gives
\[
 F(t)=\#\{k:k+a_k\le t\}+\#\{k:b_k\le t\}.
 \label{eq:counting-renewal}
\]
The second term is \(t/\beta+o(t)\).  Put
\[
 c_-=\liminf_{k\to\infty}\frac{a_k}{k},
 \qquad
 c_+=\limsup_{k\to\infty}\frac{a_k}{k}.
\]
The merge itself gives finite positive linear bounds, so \(0<c_-\le c_+<\infty\).
Monotone inversion yields
\[
 \liminf_{t\to\infty}\frac{F(t)}t=\frac1{c_+},
 \qquad
 \limsup_{t\to\infty}\frac{F(t)}t=\frac1{c_-}.
\]
The first count on the right of \eqref{eq:counting-renewal} has lower density at least
\(1/(1+c_+)\) and upper density at most \(1/(1+c_-)\).  Consequently,
\[
 \frac1{c_+}\ge\frac1{1+c_+}+\frac1\beta,
 \qquad
 \frac1{c_-}\le\frac1{1+c_-}+\frac1\beta.
\]
The function \(c\mapsto1/[c(1+c)]\) is strictly decreasing.  Both inequalities force
\(c_+\le c\le c_-\), where \(c\) is the unique solution of
\eqref{eq:merge-slope-equation}.  Hence \(c_-=c_+=c\).
\end{proof}

Since
\[
 u_k=2k+\floor{k/2}-s_2(k)=\frac52k+O(\log k),
\]
where \(s_2(k)\) is the binary digit sum, the lemma applies with \(\beta=5/2\).

\begin{corollary}[Exact two-adic linear constant]
\label{cor:p2-slope}
\begin{equation}
 \boxed{\displaystyle
 \lim_{k\to\infty}\frac{\nu_{\B,2}(k)}{k}
 =\frac{\sqrt{11}-1}{2}
 =1.1583123951777\ldots.}
 \label{eq:p2-slope}
\end{equation}
\end{corollary}

\begin{proof}
Equation \eqref{eq:merge-slope-equation} becomes
\[
 \frac1c=\frac1{1+c}+\frac25,
 \qquad
 2c(1+c)=5.
\]
The positive root is \((\sqrt{11}-1)/2\).
\end{proof}

\begin{statusbox}
The formulas in this section are \statusproved.  They use only elementary \(p\)-adic topology,
cyclic unit groups at odd primes, the order of \(3\) modulo powers of \(2\), and Bhargava's
choice-independence theorem.  They are new relative to the supplied corpus.  No claim of
priority over unpublished or unindexed work is made.
\end{statusbox}

\section{Global factorization and a finite exact algorithm}
\label{sec:global-factorial}

\subsection{A collision cutoff for every contributing prime power}

Formula~\eqref{eq:odd-prime-main} is formally infinite, but for each \(k\) only finitely many
prime powers can occur.  The following bound is deliberately elementary and proof-producing.

\begin{theorem}[Square-box cutoff]
\label{thm:cutoff}
Let \(k\ge2\), put \(R=\floor{\sqrt{k}}\), and let \(p\ge5\) be prime.  If
\(h_p(e)\le k\), then
\begin{equation}
  \boxed{p^e<6^R.}
  \label{eq:cutoff}
\end{equation}
\end{theorem}

\begin{proof}
The \((R+1)^2\) semigroup elements
\[
  2^a3^b,\qquad 0\le a,b\le R,
\]
map into a group of cardinality \(h_p(e)\le k<(R+1)^2\).  Two distinct exponent pairs
therefore collide modulo \(p^e\).  Unique factorization makes the corresponding integers
distinct, so \(p^e\) divides a nonzero difference.  Both integers lie between \(1\) and
\(6^R\); hence the absolute difference is less than \(6^R\).
\end{proof}

\begin{remark}
A simplex of exponent pairs gives slightly sharper constants, but the square cutoff has three
advantages: it is one line, it is uniform in the prime and lifting depth, and it makes exact
verification straightforward.  The bound is not intended as an asymptotic estimate for the
largest prime factor; it is a completeness certificate for the computation.
\end{remark}

\subsection{The complete factorization theorem}

\begin{theorem}[Bhargava factorial of the \(2\)--\(3\) semigroup]
\label{thm:global-factorial}
For every \(k\ge0\),
\begin{equation}
 \boxed{
 k!_{\B}
 =
 2^{a_k}\,
 3^{v_3(k!)}
 \prod_{p\ge5}
 p^{\displaystyle\sum_{e\ge1}\floor{k/h_p(e)}}.
 }
 \label{eq:global-factorial}
\end{equation}
where \(a_k\) is the exact merge sequence of
Theorem~\ref{thm:p2-recursion}.  For \(k\ge2\), every prime power that contributes to the
product satisfies \(p^e<6^{\floor{\sqrt{k}}}\).
\end{theorem}

\begin{proof}
Combine Theorems~\ref{thm:odd-prime-formula}, \ref{thm:p3},
and~\ref{thm:p2-recursion}.  The cutoff is Theorem~\ref{thm:cutoff}.
\end{proof}

This gives a finite exact procedure:

\begin{enumerate}
\item compute \(a_k\) by the two-branch merge;
\item add \(v_3(k!)\);
\item enumerate primes \(5\le p<6^R\);
\item for powers \(p^e<6^R\), compute
\[
 h_p(e)=\lcm(\ord_{p^e}(2),\ord_{p^e}(3));
\]
\item add \(\floor{k/h_p(e)}\) to the exponent of \(p\).
\end{enumerate}

No floating-point arithmetic enters.

\subsection{First exact values}

Table~\ref{tab:first-factorials} separates the structural exponents from the unramified
factor.  The last column is the exact integer gain over the ordinary factorial.

\begin{longtable}{@{}r r r l r@{}}
\caption{Exact values of \(k!_{\B}\) through \(k=20\).}
\label{tab:first-factorials}\\
\toprule
\(k\) & \(v_2(k!_{\B})\) & \(v_3(k!_{\B})\) &
\(\prod_{p\ge5}p^{v_p(k!_{\B})}\) & \(k!_{\B}/k!\)\\
\midrule
\endfirsthead
\toprule
\(k\) & \(v_2(k!_{\B})\) & \(v_3(k!_{\B})\) &
\(\prod_{p\ge5}p^{v_p(k!_{\B})}\) & \(k!_{\B}/k!\)\\
\midrule
\endhead
0 & 0 & 0 & $1$ & $1$ \\
1 & 0 & 0 & $1$ & $1$ \\
2 & 1 & 0 & $1$ & $1$ \\
3 & 1 & 1 & $1$ & $1$ \\
4 & 3 & 1 & $5$ & $5$ \\
5 & 4 & 1 & $5$ & $2$ \\
6 & 4 & 2 & $5\,7$ & $7$ \\
7 & 5 & 2 & $5\,7$ & $2$ \\
8 & 7 & 2 & $5^{2}\,7$ & $5$ \\
9 & 9 & 4 & $5^{2}\,7$ & $20$ \\
10 & 9 & 4 & $5^{2}\,7\,11$ & $22$ \\
11 & 10 & 4 & $5^{2}\,7\,11\,23$ & $92$ \\
12 & 10 & 5 & $5^{3}\,7^{2}\,11\,13\,23$ & $10465$ \\
13 & 12 & 5 & $5^{3}\,7^{2}\,11\,13\,23$ & $3220$ \\
14 & 13 & 5 & $5^{3}\,7^{2}\,11\,13\,23$ & $460$ \\
15 & 14 & 6 & $5^{3}\,7^{2}\,11\,13\,23$ & $184$ \\
16 & 15 & 6 & $5^{4}\,7^{2}\,11\,13\,17\,23$ & $1955$ \\
17 & 18 & 6 & $5^{4}\,7^{2}\,11\,13\,17\,23$ & $920$ \\
18 & 19 & 8 & $5^{4}\,7^{3}\,11\,13\,17\,19\,23$ & $122360$ \\
19 & 19 & 8 & $5^{4}\,7^{3}\,11\,13\,17\,19\,23$ & $6440$ \\
20 & 20 & 8 & $5^{6}\,7^{3}\,11^{2}\,13\,17\,19\,23$ & $177100$ \\
\bottomrule
\end{longtable}

Several features appear immediately.

\begin{itemize}
\item The first proper-orbit prime is \(23\): both \(2\) and \(3\) have order \(11\)
modulo \(23\), so \(h_{23}(1)=11\), and \(23\mid11!_{\B}\).
\item The ratio \(k!_{\B}/k!\) is an integer but need not be monotone.  A generalized
factorial is indexed by degree, not generated by multiplying successive ratios.
\item Exceptional orbit primes can be much larger than \(k\).  The prime \(431\), for
example, has \(h_{431}(1)=43\) and first enters at degree \(43\).
\end{itemize}

Selected larger values are shown in Table~\ref{tab:selected-factorials}.

\begin{table}[ht]
\centering
\small
\begin{tabular}{@{}r r r r r r r@{}}
\toprule
\(k\)&\(v_2\)&\(v_3\)&largest prime&support size&\(k!_{\B}/k!\)&
\(\log(k!_{\B}/k!)\)\\
\midrule
20 & 20 & 8 & 23 & 9 & $177100$ & 12.084470 \\
25 & 25 & 10 & 47 & 10 & $3934840$ & 15.185381 \\
30 & 31 & 14 & 47 & 12 & $412855520$ & 19.838608 \\
35 & 38 & 15 & 71 & 13 & $564887360$ & 20.152137 \\
40 & 43 & 18 & 73 & 16 & $1627306323412000$ & 35.025702 \\
45 & 48 & 21 & 431 & 18 & $1001511424282768000$ & 41.448042 \\
50 & 53 & 22 & 431 & 19 & $146685344363879921600$ & 46.434821 \\
\bottomrule
\end{tabular}
\caption{Selected exact factorial diagnostics.  ``Support size'' counts all prime divisors.}
\label{tab:selected-factorials}
\end{table}

\subsection{Independent consistency checks}

The implementation in Appendix~\ref{app:code} performs the following checks.

\begin{enumerate}
\item It verifies every reported multiplicative order by exact modular exponentiation.
\item It recomputes the \(2\)-adic stream from the residue-tree merge, rather than from a
fitted formula.
\item For small \(k\) and small primes, it compares the orbit formula with direct greedy
minimization over a large finite rectangle of exponent pairs.
\item It verifies \(k!\mid k!_{\B}\), a consequence of \(\B\subset\Z\).
\item It refuses to omit a prime power unless Theorem~\ref{thm:cutoff} places it beyond the
complete search range.
\end{enumerate}

\begin{statusbox}
The factorization theorem and the cutoff are \statusproved.  Tables~\ref{tab:first-factorials}
and~\ref{tab:selected-factorials} are \statuscomputed{} with the source printed verbatim in
Appendix~\ref{app:code}.
\end{statusbox}

\section{Local excess, global growth, and a rank-two Artin problem}
\label{sec:growth}

\subsection{Fixed-prime excess over the ordinary factorial}

For a fixed prime, floors introduce only logarithmic error.  Define
\[
 \delta_2=\frac{\sqrt{11}-1}2-1=\frac{\sqrt{11}-3}2,
 \qquad
 \delta_3=0,
\]
and for \(p\ge5\),
\begin{equation}
 \delta_p=\sum_{e\ge1}\frac1{h_p(e)}-\frac1{p-1}.
 \label{eq:delta-p}
\end{equation}

\begin{proposition}[Local excess law]
\label{prop:local-excess}
For every fixed prime \(p\),
\begin{equation}
 v_p(k!_{\B})-v_p(k!)=\delta_p k+o_p(k).
 \label{eq:local-excess-asymptotic}
\end{equation}
For \(p=3\) the error is zero, and for every fixed \(p\ge5\) it is
\(O_p(\log k)\).  Moreover,
\[
 \delta_2>0,\qquad \delta_3=0,\qquad \delta_p\ge\frac1{(p-1)^2}>0
 \quad(p\ge5).
\]
\end{proposition}

\begin{proof}
At \(p=2\), Corollary~\ref{cor:p2-slope} gives the stated \(o(k)\) remainder after
subtracting \(v_2(k!)=k+O(\log k)\).  At \(p=3\),
Theorem~\ref{thm:p3} is exact.  At \(p\ge5\), sum
\(\floor{k/h_p(e)}=k/h_p(e)+O(1)\) over the \(O_p(\log k)\) levels for which
\(h_p(e)\le k\), and compare with Legendre's formula.  Inequality
\eqref{eq:local-excess-lower} gives the last assertion.
\end{proof}

Consequently, for every finite set \(P\) of primes containing \(2\) but not \(3\),
\begin{equation}
 \log\frac{k!_{\B}}{k!}
 \ge
 k\sum_{p\in P}\delta_p\log p+o_P(k).
 \label{eq:finite-prime-lower}
\end{equation}
Already \(p=2,5,7,11,13,17,19,23\) give a certified linear coefficient exceeding
\(0.48\).  This is a genuine denominator gain, not a numerical artifact.

\subsection{Why the global asymptotic is harder}

One might expect a Stirling law
\begin{equation}
 \log k!_{\B}=k\log k+C_{\B}k+o(k).
 \label{eq:artin-bhargava-stirling}
\end{equation}
The fixed-prime analysis is compatible with such a formula, but does not prove it.  The issue
is the moving set of primes \(p>k\) for which the subgroup
\(\langle2,3\rangle\bmod p\) has unexpectedly small order.  Each such prime contributes to
\(k!_{\B}\) even though it does not divide \(k!\).

\begin{problem}[Rank-two Artin--Bhargava constant]
\label{prob:artin-bhargava}
Prove or disprove the existence of a constant \(C_{\B}\) satisfying
\eqref{eq:artin-bhargava-stirling}, and express it as a convergent combination of:
\begin{enumerate}
\item the two structural local constants at \(2\) and \(3\);
\item the orbit densities \(\sum_e1/h_p(e)\) at unramified primes;
\item a centered correction for exceptional small-order primes.
\end{enumerate}
\end{problem}

This problem is adjacent to, but not identical with, Artin's primitive-root conjecture.
Primitive-root questions ask whether a reduction has maximal order.  Here every possible
orbit size is weighted by its reciprocal, and prime-power lifting enters as well.  Average-order
theorems for finitely generated subgroups and GRH-conditional Artin-density formulas are
natural inputs, but a uniform summation strong enough for \eqref{eq:artin-bhargava-stirling}
still has to be written down.

\subsection{Unconditional growth bounds}

The ordinary factorial gives the lower bound
\begin{equation}
  \log k!\le\log k!_{\B}.
  \label{eq:growth-lower}
\end{equation}
A first upper bound comes from subset monotonicity, not from an arbitrary star product.  Let
\(D_2=\{2^n:n\in\Nzero\}\).  Since \(D_2\subset\B\), one has
\(k!_{\B}\mid k!_{D_2}\).  The natural geometric ordering is a simultaneous local ordering
of \(D_2\), so
\[
 k!_{D_2}=\prod_{j=0}^{k-1}(2^k-2^j).
\]
Consequently
\begin{equation}
  \log k!_{\B}\le k^2\log2+O(k).
  \label{eq:quadratic-upper}
\end{equation}

The rank-two geometry improves the exponent, but a cumulative-energy argument is required.
For
\[
 N(L)=\#\{(a,b)\in\Nzero^2:a\log2+b\log3\le L\},
\]
elementary lattice counting gives
\[
 N(L)=\frac{L^2}{2\log2\log3}+O(L).
\]
Write
\[
 D_m(\B)=\prod_{j=0}^{m-1}j!_{\B}.
\]
Energy divisibility, proved later in Proposition~\ref{prop:energy-divisibility}, implies that
\(D_m(\B)\) divides the Vandermonde product of any \(m\) semigroup nodes.  Also the local
sequences \(\nu_{\B,p}(k)\) are nondecreasing: the first \(k\) factors in the defining product
for stage \(k+1\) already have valuation at least \(\nu_{\B,p}(k)\).  Hence
\(k!_{\B}\mid j!_{\B}\) whenever \(j\ge k\).

\begin{proposition}[Rank-two cumulative-energy upper bound]
\label{prop:k32}
Put
\[
 C_*=\frac34\left(\frac53\right)^{5/2}
       \sqrt{2\log2\log3}
 =3.3191974798\ldots.
\]
Then
\begin{equation}
 \log k!_{\B}\le C_*k^{3/2}+O(k).
 \label{eq:k32-upper}
\end{equation}
\end{proposition}

\begin{proof}
Fix \(c>1\), let \(m=\ceil{ck}\), and choose \(m\) semigroup elements below \(e^L\), where
\[
 L=\sqrt{2\log2\log3}\,\sqrt m+O(1).
\]
Their Vandermonde product is at most \(e^{L\binom m2}\), so energy divisibility gives
\[
 \log D_m(\B)\le L\binom m2.
\]
Because every factor \(j!_{\B}\) with \(k\le j<m\) is a multiple of \(k!_{\B}\),
\[
 (m-k)\log k!_{\B}\le\log D_m(\B).
\]
It follows that
\[
 \log k!_{\B}
 \le
 \left(\frac{c^{5/2}}{2(c-1)}\sqrt{2\log2\log3}+o(1)\right)k^{3/2}.
\]
The coefficient is minimized at \(c=5/3\), giving \(C_*\).  The lattice-counting and
rounding errors contribute \(O(k)\).
\end{proof}

The use of cumulative energy is essential.  In a sparse set, \(k!_S\) need not divide the
product of the differences from an arbitrary node to \(k\) arbitrary companions; for
example, this already fails on \(\B\) at degree six.  What is universal is the full
Vandermonde divisor \(D_m(S)\).  The gap between \(k\log k\) and \(k^{3/2}\) remains
substantial.  Closing it is not required for the Alaoglu--Erd\H{o}s conjecture, but a
linear-error Stirling law would make the normalized interpolation lattice quantitatively
transparent.

\subsection{Recent determinant-completion context}

Bhargava's factorials were created to organize fixed divisors and integer-valued polynomial
bases.  Recent work studies Stirling laws for polynomial images, simultaneous \(p\)-orderings,
and determinant-like completions of generalized factorials
\cite{Takeda2023,Szumowicz2023,Diaz2026}.  In particular, a July 2026
working draft proposes renewal and determinant-line models and explicitly separates the
multivariable theory as a further problem.  The semigroup \(\B\) is a natural rank-two test:
its odd-prime pieces are procyclic orbit systems, its \(3\)-adic piece is ordinary, and its
\(2\)-adic piece is a solvable nonlinear renewal.  This report uses only the elementary local
consequences, not any conjectural analytic completion.

\begin{statusbox}
Proposition~\ref{prop:local-excess} and the growth bounds are \statusproved.
Problem~\ref{prob:artin-bhargava} is \statusopen.  Its formulation is motivated by the exact
factorial calculus, not assumed in any later argument.
\end{statusbox}

\section{A counterexample as an isomorphism of two factorial semigroups}
\label{sec:two-semigroups}

Assume for this section that \(x\) is a hypothetical nonintegral solution and put
\[
 M=2^x,\qquad N=3^x.
\]
Define the output semigroup
\[
 \Cout=\{M^aN^b:a,b\in\Nzero\}.
\]

\begin{proposition}[Power-semigroup isomorphism]
\label{prop:semigroup-isomorphism}
The map
\[
 \Phi_x:\B\longrightarrow\Cout,\qquad
 \Phi_x(2^a3^b)=M^aN^b=(2^a3^b)^x
\]
is an order-preserving monoid isomorphism.  Both \(2,3\) and \(M,N\) are multiplicatively
independent pairs.
\end{proposition}

\begin{proof}
Multiplicative independence of \(2,3\) is unique factorization.  If \(M^r=N^s\), then
\(2^{xr}=3^{xs}\); because \(x>0\), this implies \(2^r=3^s\), hence \(r=s=0\).
Thus exponent pairs give unique coordinates on both semigroups.  Multiplicativity is immediate,
and \(t\mapsto t^x\) is strictly increasing on \(\R_{>0}\).
\end{proof}

This reformulation separates two geometries.

\begin{description}
\item[Archimedean geometry.] \(\Phi_x\) preserves the total ordering and raises all ratios to
the power \(x\).
\item[Finite-place geometry.] The residue trees of \(\langle2,3\rangle\) and
\(\langle M,N\rangle\) can be completely different.  At a prime not dividing the generators,
each is controlled by the size of a generated subgroup; at a prime dividing a generator, a
ramified branch tree appears.
\end{description}

The exact input factorial is now known.  The output factorial has the same formal local
description:
for every odd prime \(p\nmid MN\),
\begin{equation}
 v_p(k!_{\Cout})
 =
 \sum_{e\ge1}
 \floor{\frac{k}{
 |\langle M,N\rangle\bmod p^e|}}.
 \label{eq:output-unramified}
\end{equation}
At primes dividing \(MN\), the closure decomposes into valuation branches, just as
\(\B\) decomposes at \(2\).  This makes the pair
\[
  \bigl(k!_{\B},k!_{\Cout}\bigr)_{k\ge0}
\]
a new invariant of a hypothetical solution.  The order isomorphism controls the real sizes of
the elements but does not, by itself, identify these two factorial sequences.

\begin{problem}[Factorial-spectrum rigidity]
\label{prob:factorial-spectrum}
Find a property of an order-preserving power isomorphism
\(\Phi_x:\B\to\Cout\) which forces a quantitative relation between
\(k!_{\B}\) and \(k!_{\Cout}\) strong enough to imply that \(x\) is an integer.
\end{problem}

A naive equality of factorials is false and should not be expected: Bhargava factorials are
arithmetic invariants of the embedded sets, not abstract monoid invariants.  Any successful
comparison has to use the special analytic form \(q\mapsto q^x\).

\section{The optimally normalized interpolation determinant}
\label{sec:normalized-determinant}

\subsection{An integer determinant forced by a counterexample}

Fix a regular basis \(F_0,F_1,\ldots\) of \(\Int(\B,\Z)\) as in
\eqref{eq:regular-basis}.  For \(d\ge0\) and distinct nodes
\[
 E=(q_0,\ldots,q_{d+1})\subset\B,
\]
define
\begin{equation}
 \Delta_{d,x}(E)
 =
 \det
 \begin{pmatrix}
 F_0(q_0)&F_1(q_0)&\cdots&F_d(q_0)&q_0^x\\
 F_0(q_1)&F_1(q_1)&\cdots&F_d(q_1)&q_1^x\\
 \vdots&\vdots&&\vdots&\vdots\\
 F_0(q_{d+1})&F_1(q_{d+1})&\cdots&F_d(q_{d+1})&q_{d+1}^x
 \end{pmatrix}.
 \label{eq:Delta-def}
\end{equation}

\begin{theorem}[Bhargava-normalized determinant]
\label{thm:normalized-determinant}
Under a hypothetical nonintegral solution \(x\),
\begin{equation}
 \boxed{\Delta_{d,x}(E)\in\Z\setminus\{0\}.}
 \label{eq:Delta-integer}
\end{equation}
Moreover,
\begin{equation}
 \Delta_{d,x}(E)
 =
 \frac{
 \det\bigl(q_i^{\lambda_j}\bigr)_{
 0\le i,j\le d+1,\,
 (\lambda_0,\ldots,\lambda_{d+1})=(0,1,\ldots,d,x)}
 }{\displaystyle\prod_{j=0}^{d}j!_{\B}}.
 \label{eq:Delta-monomial}
\end{equation}
\end{theorem}

\begin{proof}
Every \(F_j(q_i)\) is an integer, and every \(q_i^x\) is an integer by
\eqref{eq:semigroup-power}; hence the determinant is integral.  It is nonzero because the
functions
\[
 1,t,t^2,\ldots,t^d,t^x
\]
form a Chebyshev system on \(\R_{>0}\) when \(x\notin\{0,1,\ldots,d\}\).  Equivalently, after
the change \(t=e^y\), a nontrivial linear combination is an exponential polynomial with
distinct real exponents and has at most \(d+1\) real zeros.

The change of columns from \(1,T,\ldots,T^d\) to the regular basis is upper triangular with
diagonal entries \(1/j!_{\B}\).  Its determinant is
\(\prod_{j=0}^{d}(j!_{\B})^{-1}\), proving \eqref{eq:Delta-monomial}.
\end{proof}

\subsection{Divided differences and an exact inequality}

Let
\[
 V(E)=\prod_{0\le i<j\le d+1}(q_j-q_i)
\]
with the nodes listed increasingly.  The standard divided-difference identity gives
\begin{equation}
 \det(1,q_i,\ldots,q_i^d,q_i^x)
 =
 V(E)\,[q_0,\ldots,q_{d+1}]\,t^x.
 \label{eq:divided-difference-det}
\end{equation}
For some \(\xi\in(q_0,q_{d+1})\),
\begin{equation}
 [q_0,\ldots,q_{d+1}]\,t^x
 =
 \binom{x}{d+1}\xi^{x-d-1}.
 \label{eq:mean-value-divdiff}
\end{equation}
Combining this with Theorem~\ref{thm:normalized-determinant} gives a necessary inequality
for every counterexample:
\begin{equation}
 \boxed{
 \frac{|V(E)|}{\prod_{j=0}^{d}j!_{\B}}\,
 \left|\binom{x}{d+1}\right|
 \xi^{x-d-1}
 \ge1.}
 \label{eq:master-arch-ineq}
\end{equation}

This is the exact place where the generalized factorial matters.  Replacing
\(\prod j!_{\B}\) by ordinary factorials throws away genuine local divisibility; replacing it
by a node-specific product obscures which part is universal and which part is geometric.

\subsection{The factorial defect of a node set}

Applying the regular basis through degree \(m-1\) to \(m\) nodes gives:

\begin{proposition}[Energy divisibility]
\label{prop:energy-divisibility}
For every \(m\)-element set \(E\subset\B\),
\begin{equation}
  \prod_{j=0}^{m-1}j!_{\B}\mid V(E).
  \label{eq:energy-divisibility}
\end{equation}
Consequently,
\begin{equation}
  Q_{\B}(E):=
  \frac{|V(E)|}{\prod_{j=0}^{m-1}j!_{\B}}
  \in\N
  \label{eq:defect-Q}
\end{equation}
is a well-defined factorial defect.
\end{proposition}

\begin{proof}
The determinant of the integer matrix \((F_j(q_i))_{0\le i,j<m}\) is an integer.  The
triangular leading-coefficient transformation identifies that determinant with
\(V(E)/\prod_{j=0}^{m-1}j!_{\B}\), up to sign.
\end{proof}

A set with \(Q_{\B}(E)=1\) is simultaneously optimal at every prime.  Such sets are finite
prefixes of a simultaneous \(p\)-ordering.  In terms of \(Q_{\B}\),
\eqref{eq:master-arch-ineq} becomes
\begin{equation}
 Q_{\B}(E)\,(d+1)!_{\B}
 \left|\binom{x}{d+1}\right|\xi^{x-d-1}\ge1.
 \label{eq:master-Q}
\end{equation}
This isolates three independent quantities: the universal factorial, the geometric defect of
the chosen nodes, and the Archimedean curvature of \(t^x\).

\section{The first global incompatibility: an unavoidable factor \(7\)}
\label{sec:first-defect}

The first few semigroup points behave deceptively like an ordinary Newton sequence:
\[
 \{1,2\},\quad
 \{1,2,3\},\quad
 \{1,2,3,4\},\quad
 \{1,2,3,4,6\}
\]
have factorial defect \(1\).  The next step fails.

\begin{theorem}[First factorial defect]
\label{thm:first-defect}
For every six-element subset \(E\subset\B\),
\begin{equation}
  Q_{\B}(E)\ge7.
  \label{eq:first-defect-lower}
\end{equation}
Equality is attained by
\begin{equation}
  E_6=\{1,2,3,4,6,8\},
  \qquad
  Q_{\B}(E_6)=7.
  \label{eq:first-defect-equality}
\end{equation}
Thus \(\B\) has no simultaneous \(p\)-ordering whose first six points are globally optimal.
\end{theorem}

\begin{proof}
The exact denominator is
\[
 \prod_{j=0}^{5}j!_{\B}
 =1\cdot1\cdot2\cdot6\cdot120\cdot240
 =345600.
 \label{eq:six-denominator}
\]
Translate a six-point integer set so its least point is \(0\), writing
\[
 0=y_0<y_1<\cdots<y_5=L.
\]
Its Vandermonde product is translation-invariant.  The four interior points contribute at
least the ordinary four-point Vandermonde \(1!2!3!=12\), and
\(y_i(L-y_i)\ge L-1\).  Hence
\begin{equation}
 V(y_0,\ldots,y_5)
 \ge 12L(L-1)^4.
 \label{eq:span-bound-six}
\end{equation}
If \(Q_{\B}(E)<7\), then \(V(E)<7\cdot345600\), so \eqref{eq:span-bound-six} gives
\(L\le12\).  There are only \(\binom{12}{5}=792\) normalized gap patterns to inspect.

Exact enumeration leaves the following patterns below the threshold, up to reflection:
\[
\begin{array}{c|c}
\text{gap pattern}&V/345600\\ \hline
(0,1,2,3,4,5)&1/10\\
(0,1,2,3,4,6)&3/5\\
(0,1,2,3,5,6)&3/2\\
(0,1,2,4,5,6)&2\\
(0,1,2,3,4,7)&21/10\\
(0,2,3,4,5,7)&7/2\\
(0,1,2,3,4,8)&28/5
\end{array}
\]
together with reflected duplicates.  Because Proposition~\ref{prop:energy-divisibility}
requires the ratio to be an integer for a subset of \(\B\), the only possible defect below
\(7\) is \(2\), with gap pattern
\[
  (0,1,2,4,5,6).
\]
A translate of this pattern would contain two triples of consecutive \(3\)-smooth integers:
\[
 t,t+1,t+2
 \quad\hbox{and}\quad
 t+4,t+5,t+6.
\]
The only triples of consecutive elements of \(\B\) are \(1,2,3\) and \(2,3,4\).  Indeed, if
the first and third are odd, they are powers of \(3\) differing by \(2\), forcing \(1,3\).
If the middle one is odd, it is \(3^a\), while its neighbors are not divisible by \(3\) and
must be powers of \(2\); two powers of \(2\) differing by \(2\) force \(2,4\).  No translate
contains two such triples four units apart.  Therefore a defect below \(7\) is impossible.

Finally,
\[
 V(1,2,3,4,6,8)=7\cdot345600,
\]
which proves sharpness.  The finite pattern list is independently reproduced by the exact
script in Appendix~\ref{app:code}.
\end{proof}

\begin{remark}
The factor \(7\) is not inserted by hand.  The set \(E_6\) is locally optimal at \(2,3,5\);
its sole excess prime is \(7\), coming from the endpoint difference \(8-1\).  The theorem is
the first place where primewise optimal orderings cannot be synchronized globally.
\end{remark}

\begin{corollary}[No six-node Newton prefix]
\label{cor:no-six-newton}
There is no sequence \(a_0,\ldots,a_5\in\B\) which is a \(p\)-ordering prefix for every prime
\(p\) simultaneously.
\end{corollary}

\begin{warning}
The absence of a simultaneous ordering does not prevent a regular basis of
\(\Int(\B,\Z)\).  The regular basis is assembled from compatible characteristic ideals, not
from one universal Newton sequence.  Confusing these two facts would incorrectly discard the
Bhargava lattice just when its first nontrivial global obstruction appears.
\end{warning}

\begin{statusbox}
The normalized determinant, energy divisibility, and Theorem~\ref{thm:first-defect} are
\statusproved.  The 792-pattern check is finite exact arithmetic; the source and threshold
proof are included, so no heuristic search boundary is involved.
\end{statusbox}

\section{Subunit locking in the exact Bhargava lattice}
\label{sec:locking}

\subsection{The optimal integer-valued approximation statement}

\begin{theorem}[Bhargava locking]
\label{thm:bhargava-locking}
Assume \(x\notin\Z\) and \(q^x\in\Z\) for every \(q\in\B\).  Let \(d\ge0\) and let
\(q_0,\ldots,q_{d+1}\) be distinct elements of \(\B\).  Then
\begin{equation}
 \inf_{P\in\Int_{\le d}(\B,\Z)}
 \max_{0\le i\le d+1}|q_i^x-P(q_i)|
 \ge1.
 \label{eq:locking-distance}
\end{equation}
\end{theorem}

\begin{proof}
For \(P\in\Int(\B,\Z)\), every residual \(q_i^x-P(q_i)\) is an integer.  If all residuals
had absolute value below \(1\), they would all vanish.  The nonzero generalized polynomial
\(t^x-P(t)\) would then have \(d+2\) positive zeros, contradicting the Chebyshev property of
\(1,t,\ldots,t^d,t^x\).
\end{proof}

This is an exact closest-vector obstruction.  In a regular basis write
\[
 P=\sum_{j=0}^{d}c_jF_j,\qquad c_j\in\Z.
\]
For \(E=(q_0,\ldots,q_{d+1})\), let
\[
 A_E=(F_j(q_i))\in\Z^{(d+2)\times(d+1)},
 \qquad
 y_x(E)=(q_i^x)\in\Z^{d+2}.
\]
Then Theorem~\ref{thm:bhargava-locking} is
\begin{equation}
 \dist_{\ell^\infty}\bigl(y_x(E),A_E\Z^{d+1}\bigr)\ge1.
 \label{eq:CVP}
\end{equation}
The matrix \(A_E\) is the optimal integer-valued-polynomial lattice; no larger universal
denominators are available.

\subsection{A determinant dual certificate}

The left kernel of \(A_E\) has rank one.  Let
\(\lambda(E)\in\Z^{d+2}\) be its primitive generator, with alternating signs for increasing
nodes.  Then
\begin{equation}
 \lambda(E)^{\!T}y_x(E)\in\Z\setminus\{0\},
 \qquad
 |\lambda(E)^{\!T}y_x(E)|\ge1.
 \label{eq:dual-certificate}
\end{equation}
This is the cofactor form of Theorem~\ref{thm:normalized-determinant}.  It turns the search for
a contradiction into a finite task: construct \(E\) for which the exact integer linear form
has a rigorous Archimedean upper bound below one.

\subsection{Four concrete proof criteria}

\begin{criterion}[Normalized determinant criterion]
\label{crit:det}
Conjecture~\ref{conj:AE} follows if, for every sufficiently large possible counterexample
\(x\), one can find \(d\) and \(E\subset\B\), \(|E|=d+2\), such that
\begin{equation}
 \left|
 \det(1,q_i,\ldots,q_i^d,q_i^x)
 \right|
 <
 \prod_{j=0}^{d}j!_{\B}.
 \label{eq:det-criterion}
\end{equation}
\end{criterion}

\begin{proof}
The normalized determinant would be a nonzero integer of absolute value below one.
\end{proof}

\begin{criterion}[Low-defect curvature criterion]
\label{crit:low-defect}
It is enough to find \(d,E\), and a valid divided-difference bound such that
\begin{equation}
 Q_{\B}(E)(d+1)!_{\B}
 \left|\binom{x}{d+1}\right|
 \max_{t\in[\min E,\max E]}t^{x-d-1}<1.
 \label{eq:low-defect-criterion}
\end{equation}
\end{criterion}

\begin{criterion}[Regular-basis closest vector]
\label{crit:cvp}
It is enough to exhibit \(P\in\Int_{\le d}(\B,\Z)\) and \(d+2\) semigroup nodes on which
\[
 |q^x-P(q)|<1.
\]
The coefficients of \(P\) may be searched directly in the regular basis; no interpolation
node set is required.
\end{criterion}

\begin{criterion}[Output-input paired lattice]
\label{crit:paired}
Let \(G_0,G_1,\ldots\) be a regular basis for
\(\Int(\Cout,\Z)\).  It is enough to combine the input determinant for \(t^x\) and the output
determinant for \(t^{1/x}\) into a nonzero rational number whose valuations are nonnegative at
every finite prime and whose Archimedean absolute value is below one.
\end{criterion}

The fourth criterion is deliberately phrased adelically.  Separate real estimates for a
function and its inverse do not automatically multiply to something small; the finite-place
normalizations must be coordinated.  This is a precise cancellation problem, not a missing
factorial calculation.

\subsection{Why the new lattice is useful even before it closes}

The exact factorial performs three reductions.

\begin{enumerate}
\item It replaces a vague search over rational polynomial coefficients by the canonical
lattice \(\Int_{\le d}(\B,\Z)\).
\item It identifies the maximal automatic denominator prime by prime, so any further gain must
come from cancellation, special node geometry, or the output values.
\item It supplies short finite certificates: a factorization of \(j!_{\B}\), an integer
evaluation matrix, and one determinant or dual cofactor vector.
\end{enumerate}

The same facts also explain the limit.  Generalized factorials are phase-blind.  They know how
many previous nodes occupy a residue class, but not the signed cancellation among determinant
terms.  A proof must add information not contained in the local minima alone.

\section{Localizing the semigroup: dense Archimedean orbits and a finite-prime tax}
\label{sec:localized}

\subsection{The full \(2\)--\(3\) group over \(\Z[1/6]\)}

The positive group
\[
 \GammaTwoThree=\{2^a3^b:a,b\in\Z\}
\]
is dense in \(\R_{>0}\), because \(a\log2+b\log3\) is dense in \(\R\).  It is naturally a
subset of the Dedekind domain
\[
 R_6=\Z[1/6].
\]
The nonzero prime ideals of \(R_6\) correspond to rational primes \(p\ge5\).  Bhargava's
construction therefore gives a localized factorial
\begin{equation}
 k!_{\GammaTwoThree,R_6}
 =
 \prod_{p\ge5}
 p^{\displaystyle\sum_{e\ge1}\floor{k/h_p(e)}}
 \quad\text{up to a unit of }R_6.
 \label{eq:localized-factorial}
\end{equation}
The proof is exactly Theorem~\ref{thm:odd-prime-formula}: negative exponents do not change the
finite image, and the structural primes are inverted.

This version is conceptually clean.  All compact unramified orbit data are retained, while the
two noncompact valuation directions are isolated as units.

\subsection{What a counterexample gives after localization}

Under a hypothetical solution,
\[
 \Phi_x(\GammaTwoThree)
 =
 \Gamma_{M,N}:=\{M^aN^b:a,b\in\Z\}
 \subset \Z[1/(MN)].
\]
Let
\[
 A=\Z[1/(6MN)].
\]
Choose a regular basis \(F_0^\times,F_1^\times,\ldots\) of
\(\Int(\GammaTwoThree,R_6)\).  For distinct \(q_0,\ldots,q_{d+1}\in\GammaTwoThree\), define
\[
 \Delta^\times_{d,x}(E)
 =
 \det\bigl(F_0^\times(q_i),\ldots,F_d^\times(q_i),q_i^x\bigr).
\]
The same Chebyshev argument gives
\begin{equation}
 \Delta^\times_{d,x}(E)\in A\setminus\{0\}.
 \label{eq:localized-delta}
\end{equation}
At every prime \(p\nmid6MN\), its \(p\)-adic valuation is nonnegative.  Thus every possible
denominator is supported on the finite set
\[
 S=\{p:p\mid6MN\}.
\]

For a nonzero rational \(z\) integral outside \(S\), define its denominator budget
\[
 \mathfrak d_S(z)=\prod_{p\in S}p^{\max(0,-v_p(z))}.
\]
Then
\begin{equation}
 |z|_\infty\ge\frac1{\mathfrak d_S(z)}.
 \label{eq:S-lower}
\end{equation}
Consequently:

\begin{criterion}[Finite-support compact-orbit criterion]
\label{crit:finite-support}
Conjecture~\ref{conj:AE} follows if one can construct localized node sets \(E\) for which
\begin{equation}
 |\Delta^\times_{d,x}(E)|_\infty\,
 \mathfrak d_S(\Delta^\times_{d,x}(E))<1.
 \label{eq:finite-support-criterion}
\end{equation}
\end{criterion}

All primes outside \(S\) are already handled optimally by the generalized factorial.  The
criterion asks only for output-sensitive control at finitely many structural primes.

\subsection{Why density alone loses}

Take a convergent \(a/b\) to \(\log_2 3\) and put
\[
 u=\frac{2^a}{3^b}.
\]
Then \(u\to1\), and for fixed \(d\) the geometric cluster
\[
 E_u=\{1,u,u^2,\ldots,u^{d+1}\}
\]
has
\begin{equation}
 |V(E_u)|\asymp_d|\log u|^{\binom{d+2}{2}}.
 \label{eq:cluster-vandermonde}
\end{equation}
Continued fractions give \(|\log u|\ll1/b\), so the Archimedean determinant decays only
polynomially in the exponent height.

The structural denominator grows exponentially.  In the monomial determinant, the entry
\((u^i)^j\) has a \(3\)-power denominator at most \(3^{bij}\), and
\[
 (u^i)^x=\left(\frac{M^a}{N^b}\right)^i
\]
has an \(N\)-power denominator \(N^{bi}\).  Clearing rows gives the explicit upper budget
\begin{equation}
 \mathfrak d_S
 \le
 3^{bd\binom{d+2}{2}}\,
 N^{b\binom{d+2}{2}}
 \times(\text{basis-unit factors}).
 \label{eq:row-clear-budget}
\end{equation}
For fixed \(d\), this is exponential in \(b\), while
\eqref{eq:cluster-vandermonde} is a negative power of \(b\).  Therefore the direct
``choose many group elements close to \(1\)'' strategy cannot work with termwise denominator
clearing.

\begin{proposition}[One-ratio cluster no-go]
\label{prop:cluster-nogo}
Any proof based only on:
\begin{enumerate}
\item a single ratio \(u=2^a3^{-b}\to1\);
\item the geometric nodes \(1,u,\ldots,u^{d+1}\) with fixed \(d\);
\item rowwise or termwise clearing at primes dividing \(6MN\);
\end{enumerate}
has an exponential denominator budget and only polynomial Archimedean decay.  Those
termwise estimates cannot certify \eqref{eq:finite-support-criterion} for large \(b\).
\end{proposition}

The proposition does not rule out cancellation in the determinant.  On the contrary, it
identifies cancellation as the only surviving mechanism in this localized framework.

\subsection{The finite-place cancellation target}

For a chosen cluster, expand
\[
 \Delta^\times_{d,x}(E)=\sum_{\sigma}\sgn(\sigma)\,T_\sigma
\]
into determinant terms.  The automatic valuation lower bound is the minimum of the
\(v_p(T_\sigma)\).  A proof needs a strict gain:
\begin{equation}
 v_p\!\left(\sum_\sigma\sgn(\sigma)T_\sigma\right)
 >
 \min_\sigma v_p(T_\sigma)
 \label{eq:cancellation-gain}
\end{equation}
at enough primes in \(S\), or a structural reason that the apparent denominator cancels before
the determinant is evaluated.

This is a finite-prime, output-sensitive target.  The input factorial has already extracted
all universal unramified divisibility, and Proposition~\ref{prop:cluster-nogo} has removed
termwise structural clearing.  What remains is not another counting estimate; it is the
arithmetic of the initial forms of one specific determinant.

\begin{statusbox}
The localized factorial, finite-support criterion, and one-ratio no-go are \statusproved.
The cancellation inequality \eqref{eq:cancellation-gain} is \statusopen{} and is one of the
highest-priority bridges identified in this report.
\end{statusbox}

\section{A general local calculus for two-generator semigroups}
\label{sec:general-pair}

The input computation extends to any pair of positive integers.  This is needed because the
output generators \(M,N\) are unknown.

Let
\[
 S(A,B)=\{A^aB^b:a,b\in\Nzero\}.
\]

\subsection{Unramified odd primes}

\begin{theorem}[General unramified orbit formula]
\label{thm:general-unramified}
Let \(p\) be odd and \(p\nmid AB\).  Put
\[
 h_{A,B;p}(e)
 =
 |\langle A,B\rangle\subset(\Z/p^e\Z)^\times|.
\]
Then
\begin{equation}
 v_p(k!_{S(A,B)})
 =
 \sum_{e\ge1}\floor{\frac{k}{h_{A,B;p}(e)}}.
 \label{eq:general-unramified}
\end{equation}
\end{theorem}

\begin{proof}
The proof of Theorem~\ref{thm:odd-prime-formula} uses only that the two generators are
\(p\)-adic units and that \(\Z_p^\times\) is procyclic.
\end{proof}

\subsection{Ramified primes as rooted merge systems}

When \(p\mid AB\), partition the closure by exact \(p\)-adic valuation.  Distinct valuation
branches have a predictable cross-difference valuation:
\[
 v_p(x-y)=\min(v_p(x),v_p(y))
 \quad\text{when }v_p(x)\ne v_p(y).
\]
Within a fixed branch, divide by the common power of \(p\); what remains is a compact unit
orbit or a rescaled copy of the whole closure.  The global \(p\)-ordering is therefore a
greedy merge of local cost streams.

For \((A,B)=(2,3)\) at \(p=2\), there are exactly two self-similar branches and the merge is
Theorem~\ref{thm:p2-recursion}.  For a general output pair, there may be several branches or both generators may be
ramified.  When exactly one generator is ramified, Section~\ref{sec:one-ramified} proves a
closed two-stream recursion.  When both are ramified, overlap between the scaled branches is
the unresolved finite-state problem isolated below.

\begin{problem}[Output-prime tomography]
\label{prob:output-tomography}
For a hypothetical pair \((M,N)\), compute the ramified merge systems at every
\(p\mid MN\), and relate their renewal slopes to:
\[
 v_p(M),\quad v_p(N),\quad
 \overline{\langle M p^{-v_p(M)},N p^{-v_p(N)}\rangle}\subset\Z_p^\times.
\]
Determine whether the collection of these slopes is compatible with an order-preserving power
map from \(\B\).
\end{problem}

This is a finite calculation once \(M,N\) are specified.  A proof would need a uniform
inequality valid for every possible output pair, or a descent reducing the possible valuation
profiles.

\subsection{A useful separation of roles}

The two-generator local calculus separates three ingredients that are often conflated.

\begin{description}
\item[Orbit size.] At unramified primes, this is \(h_{A,B;p}(e)\).
\item[Valuation branching.] At ramified primes, this records how many scaled copies or unit
orbits occur at each depth.
\item[Global synchronization.] A regular basis exists primewise, but a single sequence of
nodes may fail to be optimal at all primes; Theorem~\ref{thm:first-defect} detects the first
failure for \((2,3)\).
\end{description}

Any proposed determinant proof can now be audited against these three layers.

\section{The first regular-basis element that is not a Newton product}
\label{sec:first-nonnewton}

The failure of a simultaneous ordering at six nodes does not destroy the regular basis.  It
forces the local optimal polynomials to be glued coefficientwise.  This section makes that
repair completely explicit.

For degrees at most five one may use the nested node sets
\[
 \varnothing,
 \quad\{1\},
 \quad\{1,2\},
 \quad\{1,2,3\},
 \quad\{1,2,3,4\},
 \quad\{1,2,3,4,6\}.
\]
Their successive factorial defects are one.  Hence the polynomials
\begin{equation}
 F_k^{\mathrm N}(T)
 =\frac{\prod_{a\in E_k}(T-a)}{k!_{\B}}
 \qquad(0\le k\le5)
 \label{eq:newton-through-five}
\end{equation}
are integer-valued on \(\B\) and have the optimal leading coefficients.
At degree six no six-element node set supplies such a global Newton numerator, by
Theorem~\ref{thm:first-defect}.

\begin{theorem}[An explicit degree-six regular-basis element]
\label{thm:explicit-F6}
Put
\begin{align}
 P_6(T)={}&T^6+840T^5+175T^4-1260T^3 \\
          &\quad-2156T^2+1680T+720.
 \label{eq:P6-explicit}
\end{align}
Then
\begin{equation}
 F_6(T)=\frac{P_6(T)}{5040}\in\Int(\B,\Z),
 \qquad
 \operatorname{lc}(F_6)=\frac1{6!_{\B}}.
 \label{eq:F6-explicit}
\end{equation}
Consequently
\[
 1,F_1^{\mathrm N},\ldots,F_5^{\mathrm N},F_6
\]
is a regular basis through degree six.
\end{theorem}

\begin{proof}
The denominator is
\[
 6!_{\B}=2^4\,3^2\,5\,7=5040.
\]
The polynomial in~\eqref{eq:P6-explicit} is the coefficientwise Chinese-remainder gluing of
four monic local numerators:
\begin{align}
 P_6(T)&\equiv (T-1)(T-2)(T-3)(T-4)(T-6)(T-8) &&\pmod {16},
 \label{eq:P6-mod16}\\
 P_6(T)&\equiv T(T-1)(T-2)(T-3)(T-4)(T-5) &&\pmod 9,
 \label{eq:P6-mod9}\\
 P_6(T)&\equiv T^2(T^4-1) &&\pmod 5,
 \label{eq:P6-mod5}\\
 P_6(T)&\equiv T^6-1 &&\pmod 7.
 \label{eq:P6-mod7}
\end{align}
For \(q\in\B\), the possible residues modulo \(16\) are
\[
 0,1,2,3,4,6,8,9,11,12;
\]
direct substitution in the product in~\eqref{eq:P6-mod16} gives zero modulo \(16\).
The product of six consecutive integers in~\eqref{eq:P6-mod9} is divisible by \(6!\), hence
by \(9\).  At the prime \(5\), every \(q\in\B\) is a unit and satisfies \(q^4\equiv1\pmod5\).
At the prime \(7\), it is a unit and satisfies \(q^6\equiv1\pmod7\).  Thus every prime-power
factor of \(5040\) divides \(P_6(q)\), so \(F_6(q)\in\Z\).
Its leading coefficient is \(1/5040=1/6!_{\B}\), which is optimal by the characteristic-ideal
theorem.  Adjoining it to any regular basis through degree five gives the asserted basis.
\end{proof}

\begin{remark}[What the factor seven actually means]
The obstruction in Theorem~\ref{thm:first-defect} is not a missing denominator.  The optimal
denominator \(5040\) exists and the polynomial~\eqref{eq:P6-explicit} realizes it.  What fails
is the stronger demand that one monic numerator factor into six linear terms with all six
roots in \(\B\).  The first global incompatibility is therefore a coefficient-gluing defect,
not a failure of integer-valued interpolation.
\end{remark}

\begin{statusbox}
Theorem~\ref{thm:explicit-F6} is \statusproved.  Its four congruences and the displayed CRT
coefficients are also checked by the exact script in Appendix~\ref{app:code}.
\end{statusbox}

\section{Optimal factorial energies through eight nodes}
\label{sec:energies-eight}

The first-defect theorem can be pushed two steps farther.  This produces a sharp finite
picture of how quickly global synchronization deteriorates.

For \(m\ge1\), define
\begin{equation}
 \mu_m(\B)=
 \min_{\substack{E\subset\B\\|E|=m}}Q_{\B}(E).
 \label{eq:mu-m}
\end{equation}
The minimum exists because the Vandermonde product tends to infinity with the diameter of an
integer set.

\subsection{A general span bound}

\begin{lemma}[Integer-span energy bound]
\label{lem:span-general}
Let \(0=y_0<y_1<\cdots<y_{m-1}=L\) be distinct integers, with \(m\ge3\).  Then
\begin{equation}
 V(y_0,\ldots,y_{m-1})
 \ge
 \left(\prod_{j=1}^{m-3}j!\right)L(L-1)^{m-2}.
 \label{eq:span-general}
\end{equation}
\end{lemma}

\begin{proof}
The Vandermonde product of the \(m-2\) interior integers is minimized by consecutive integers
and is therefore at least \(\prod_{j=1}^{m-3}j!\).  The endpoint pair contributes \(L\), and
for every interior point \(y_i\) one has \(y_i(L-y_i)\ge L-1\).  Multiplication gives the
claim.
\end{proof}

\subsection{The exact minima}

\begin{theorem}[First eight optimal energies]
\label{thm:first-eight-energies}
The exact defect minima are
\begin{equation}
 \boxed{
 (\mu_1,\mu_2,\ldots,\mu_8)
 =(1,1,1,1,1,7,7,396).}
 \label{eq:first-eight-mu}
\end{equation}
For \(m=6,7,8\), the minimizers are unique and equal respectively to
\begin{align}
 E_6&=\{1,2,3,4,6,8\},
 \label{eq:E6-unique}\\
 E_7&=\{1,2,3,4,6,8,9\},
 \label{eq:E7-unique}\\
 E_8&=\{1,2,3,4,6,8,9,12\}.
 \label{eq:E8-unique}
\end{align}
Their defects are
\[
 Q_{\B}(E_6)=7,
 \qquad
 Q_{\B}(E_7)=7,
 \qquad
 Q_{\B}(E_8)=396.
\]
\end{theorem}

\begin{proof}
The values through \(m=5\) follow from the Newton prefixes in
\eqref{eq:newton-through-five}.  Theorem~\ref{thm:first-defect} proves the value at six.  To
see uniqueness, the equality equation
\[
 V(E)=7\prod_{j=0}^{5}j!_{\B}
\]
and Lemma~\ref{lem:span-general} restrict the normalized gap pattern to
\[
 (0,1,2,3,5,7)
 \quad\text{or its reflection}\quad
 (0,2,4,5,6,7).
\]
The first contains four consecutive integers.  The only four consecutive elements of \(\B\)
are \(1,2,3,4\): this follows immediately from the classification of consecutive triples in
the proof of Theorem~\ref{thm:first-defect}.  Hence its translate is exactly \(E_6\).  The
reflected pattern would place its four-consecutive run at a positive offset and has no positive
translate in \(\B\).

For seven nodes,
\[
 D_7:=\prod_{j=0}^{6}j!_{\B}=1741824000.
\]
If \(Q_{\B}(E)<7\), Lemma~\ref{lem:span-general} gives \(L\le19\).  Exhausting the
\(\binom{19}{6}=27132\) normalized patterns leaves, up to reflection, only
\[
\begin{array}{c|c}
 Q&\text{pattern}\ \\ \hline
 2&(0,1,3,4,5,6,8)\\
 3&(0,1,2,3,5,6,8)\\
 3&(0,1,2,3,4,5,10).
\end{array}
\]
Every pattern contains four consecutive integers; forcing that run to be \(1,2,3,4\) makes
one of the remaining entries a non-smooth integer.  Thus \(Q\ge7\).  Equality leaves four
normalized patterns, paired by reflection:
\[
\begin{split}
 &(0,1,2,3,5,7,8),\quad(0,1,2,3,6,7,8),\\
 &(0,1,2,5,6,7,8),\quad(0,1,3,5,6,7,8).
\end{split}
\]
Only the first has a positive smooth translate, namely \(E_7\).

For eight nodes,
\[
 D_8:=\prod_{j=0}^{7}j!_{\B}=17557585920000.
\]
The set \(E_8\) has
\[
 V(E_8)=396D_8.
\]
If a smaller defect existed, Lemma~\ref{lem:span-general} would force \(L\le42\).  There are
exactly
\[
 \sum_{L=7}^{42}\binom{L-1}{6}=\binom{42}{7}=26978328
\]
normalized integer patterns.  Exact integer filtering by
\(0<V(E)/D_8<396\) leaves only \(92\), with the defect histogram printed in
Table~\ref{tab:defect-eight-histogram}.  Of these, \(74\) contain four consecutive offsets.
A positive smooth translate is then possible only when the run starts at offset zero and the
translation is one; direct evaluation of the remaining offsets eliminates all \(17\) such
patterns.  The other \(18\) patterns contain a consecutive triple.  Its translated starting value must
be one or two.  Thus, for a triple beginning at normalized offset \(r\), the only candidate
translations are \(t=1-r\) and \(t=2-r\), retained only when \(t>0\); in particular only
\(r=0\) or \(r=1\) can contribute.  Every such candidate is checked exactly and fails.
Therefore no defect below \(396\) occurs.

At equality, exactly eight normalized patterns survive the Vandermonde equation.  The same
consecutive-run test eliminates seven; the sole smooth translate is
\[
 1+(0,1,2,3,5,7,8,11)=E_8.
\]
The complete enumeration and all membership checks are reproduced by the verifier in
Appendix~\ref{app:defect-code}.
\end{proof}

\begin{table}[ht]
\centering
\small
\begin{tabular}{@{}r r@{\qquad}r r@{\qquad}r r@{}}
\toprule
Defect&Count&Defect&Count&Defect&Count\\
\midrule
2&2&3&2&12&2\\
15&2&18&6&20&2\\
22&2&30&2&33&4\\
60&4&63&2&77&2\\
78&2&84&2&99&2\\
105&2&108&4&110&4\\
128&2&132&4&162&2\\
180&2&198&2&220&2\\
240&2&252&6&270&4\\
286&4&288&2&297&2\\
308&2&324&2&330&4\\
378&2&&&\\
\bottomrule
\end{tabular}
\caption{Exact histogram of normalized eight-point integer patterns with defect below
\(396\).  None has a translate contained in \(\B\).}
\label{tab:defect-eight-histogram}
\end{table}

\begin{corollary}[The synchronization collapse is abrupt]
\label{cor:synchronization-collapse}
A globally optimal node set exists through five nodes, incurs only one factor \(7\) at six and
seven nodes, and then pays at least
\[
 396=2^2\cdot3^2\cdot11
\]
at eight nodes.  The defect is therefore not governed by a gradual loss at one fixed prime.
It is a genuinely global compatibility cost.
\end{corollary}

\begin{statusbox}
Lemma~\ref{lem:span-general} is \statusproved.  Theorem~\ref{thm:first-eight-energies} combines
that proof with \statuscomputed{} finite certificates.  The search is exhaustive because the
span bounds are proved before enumeration; no empirical height cutoff is used.
\end{statusbox}

\section{A closed ramified recursion for one structural generator}
\label{sec:one-ramified}

The two-adic recursion is not peculiar to \((2,3)\).  It is the universal local model when
exactly one generator is divisible by the prime.

\begin{theorem}[One-ramified-generator merge]
\label{thm:one-ramified-merge}
Let \(A,B>1\), let \(p\mid A\) and \(p\nmid B\), and write
\[
 A=p^r u,
 \qquad r=v_p(A)\ge1,
 \qquad p\nmid u.
\]
Let
\[
 S=S(A,B),
 \qquad
 H=\overline{\langle B\rangle}^{\,p}\subset\Z_p^\times,
\]
and put
\[
 a_k=\nu_{S,p}(k),
 \qquad
 b_k=\nu_{H,p}(k).
\]
Then, as ordered multisets,
\begin{equation}
 \boxed{
 \{a_k:k\ge0\}
 =
 \{rj+a_j:j\ge0\}\uplus\{b_j:j\ge0\}.}
 \label{eq:general-ramified-merge}
\end{equation}
Equivalently, the \(p\)-ordering is obtained by greedily merging a scaled-copy stream
\(rj+a_j\) and the unit-orbit stream \(b_j\).
\end{theorem}

\begin{proof}
Let \(K=\overline S^{\,p}\).  Splitting according to whether the exponent of \(A\) is zero
gives the disjoint decomposition
\begin{equation}
 K=H\;\sqcup\;p^ruK.
 \label{eq:general-K-split}
\end{equation}
Every cross-branch difference is a unit, hence contributes zero valuation.  With \(j\)
previous points in the second branch, multiplication by \(p^ru\) contributes \(rj\) and
reduces the internal problem to \(K\), giving cost \(rj+a_j\).  With \(j\) previous unit
points, the cost is \(b_j\).  The global greedy rule is their ordered merge.
\end{proof}

\begin{corollary}[Closed renewal slope]
\label{cor:general-ramified-slope}
Suppose \(b_k/k\to\beta>0\).  Then
\begin{equation}
 \boxed{
 \lim_{k\to\infty}\frac{a_k}{k}
 =c(r,\beta)
 :=\frac{\sqrt{r^2+4r\beta}-r}{2}.}
 \label{eq:general-ramified-slope}
\end{equation}
\end{corollary}

\begin{proof}
The counting-function proof of Lemma~\ref{lem:merge-slope} applies with the first stream
having slope \(r+c\).  Thus
\[
 \frac1c=\frac1{r+c}+\frac1\beta.
\]
This is equivalent to \(c(r+c)=r\beta\), whose positive root is
\eqref{eq:general-ramified-slope}.
\end{proof}

For the input pair at \(p=2\), one has \(r=1\) and \(\beta=5/2\), recovering
Corollary~\ref{cor:p2-slope}.  For a hypothetical output pair \((M,N)\), every prime with
\(p\mid M\) and \(p\nmid N\) has the explicit factorial slope
\begin{equation}
 \lim_{k\to\infty}\frac{v_p(k!_{\Cout})}{k}
 =c\!\left(v_p(M),\sum_{e\ge1}\frac1{\ord_{p^e}(N)}\right)
 \label{eq:output-asymmetric-slope}
\end{equation}
for odd \(p\), with the evident closed-orbit interpretation at \(p=2\).

\begin{remark}
The unit factor \(u=A/p^r\) disappears from the renewal law.  It rotates the scaled copy by a
unit but does not change any difference valuation.  Thus an asymmetric support prime is
completely controlled by one ramification index and one unit-orbit density.
\end{remark}

\begin{problem}[Shared-support renewal]
\label{prob:shared-support-renewal}
When \(p\mid A\) and \(p\mid B\), the two scaled copies overlap and
\eqref{eq:general-K-split} is no longer disjoint.  Construct a finite min-plus renewal system
from the numerical semigroup generated by \(v_p(A)\) and \(v_p(B)\), including the unit cosets,
and extract its asymptotic slope and first correction term.
\end{problem}

\begin{statusbox}
The merge theorem and the closed slope are \statusproved.  They turn the previously qualitative
``output-prime tomography'' proposal into an exact formula at every prime dividing exactly one
output generator.  The shared-support case remains \statusopen.
\end{statusbox}

\section{What the new factorial calculus changes in the proof search}
\label{sec:what-changes}

The preceding results do not prove Conjecture~\ref{conj:AE}.  They do, however, change the
shape of several proof attempts in a precise way.

\subsection{What is now exact}

A hypothetical counterexample supplies an integral column \((q_i^x)\) on every input node set.
The optimal universal denominator of the polynomial columns is no longer an unknown Smith
factor: it is
\[
 \prod_{j=0}^{d}j!_{\B},
\]
with every local exponent explicit.  The first non-Newton basis element is explicit, and the
best global node energies are known through eight points.  At any asymmetric output prime, the
output factorial has the closed renewal slope~\eqref{eq:output-asymmetric-slope}.

These facts remove three ambiguities from determinant experiments:
\begin{enumerate}
\item a small determinant may no longer be credited to a denominator that is merely an artifact
of the monomial basis;
\item a node set may be scored by its exact defect \(Q_{\B}(E)\), rather than by raw pairwise
gaps;
\item output-prime gains can be separated into automatic renewal mass and genuine signed
cancellation.
\end{enumerate}

\subsection{What remains genuinely missing}

After optimal normalization, the integer determinant still contains the factor
\[
 \binom{x}{d+1}\xi^{x-d-1}.
\]
For the known counterexample-free range, a hypothetical \(x\) is already large.  Fixed-degree
curvature therefore grows rather than shrinks on positive integer nodes.  Increasing the degree
introduces more factorial mass, but Theorem~\ref{thm:first-eight-energies} shows that globally
compatible node sets pay an additional arithmetic energy that rapidly ceases to be one.

Localization removes the two input structural primes from the factorial, but
Proposition~\ref{prop:cluster-nogo} shows that a one-ratio cluster then pays an exponential
output denominator.  The only unexhausted mechanism is cancellation in the structural-prime
initial forms of the normalized determinant.  The factorial calculus identifies this as the
frontier instead of another undifferentiated ``try a larger determinant'' instruction.

\subsection{A sharpened conditional criterion}

For a localized determinant \(\Delta^\times\), write at each structural prime
\[
 \Delta^\times=p^{m_p}\bigl(I_p+pJ_p\bigr),
\]
where \(m_p\) is the minimum term valuation and \(I_p\) is the signed initial form over
\(\F_p\).  The automatic factorial and row-clearing contributions determine \(m_p\).
The missing gain is exactly the vanishing of \(I_p\).

\begin{criterion}[Initial-form bridge]
\label{crit:initial-form-bridge}
Conjecture~\ref{conj:AE} follows if one can choose a family of localized regular-basis
node sets for which
\begin{enumerate}
\item the Archimedean determinant has logarithm \(-A_d+o(A_d)\) with \(A_d\to\infty\);
\item at primes \(p\mid6MN\), the initial forms vanish to total weighted depth exceeding the
termwise denominator budget by \(A_d/\log p+o(A_d)\);
\item the determinant is nonzero, for example by total positivity or a Chebyshev argument.
\end{enumerate}
\end{criterion}

This is a sufficient condition, not a claim that such a family exists.  Its advantage is that
every universal local contribution has already been subtracted.  Only the output-specific
initial forms remain.

\section{A prioritized research program}
\label{sec:research-program}

\subsection{Priority A: structural-prime initial forms}

For small degrees, build the explicit regular basis by local CRT gluing as in
Theorem~\ref{thm:explicit-F6}.  Evaluate the localized determinant symbolically in the
valuation variables
\[
 v_p(M),\quad v_p(N),\quad M p^{-v_p(M)},\quad N p^{-v_p(N)}.
\]
The first target is not a large valuation bound.  It is a classification of when the initial
form \(I_p\) vanishes.  A useful theorem would say that simultaneous nonvanishing at all
structural primes is incompatible with the real ordering slope
\(\log M/\log2=\log N/\log3\).

\subsection{Priority B: paired regular-basis determinants}

Construct an input determinant for \(t^x\) and an output determinant for \(t^{1/x}\), each in
its own Bhargava regular basis.  Their Archimedean divided differences have reciprocal
curvatures, while their finite-place normalizations are governed by the two factorial spectra.
The concrete target is an adelic product in which the leading structural denominators cancel
before absolute values are taken.  Separate inequalities for the two determinants are too
weak.

\subsection{Priority C: low-defect node classification}

The exact sequence
\[
 1,1,1,1,1,7,7,396
\]
invites a systematic study of \(\mu_m(\B)\).  The next useful outputs are:
\begin{enumerate}
\item a lower bound \(\log\mu_m\gg m^2\) or a counterexample to it;
\item a structural description of minimizers in exponent-lattice coordinates;
\item a comparison between minimizers of \(Q_{\B}\) and minimizers of the full determinant
including the curvature factor.
\end{enumerate}
The finite verifier already supports branch-and-bound extensions beyond eight nodes, but an
unproved search cutoff must not be confused with a theorem.

\subsection{Priority D: the global Stirling law}

Resolve Problem~\ref{prob:artin-bhargava}.  Even a conditional GRH theorem with a power-saving
error would quantify how much denominator is available at moving exceptional-order primes.
The result would be independently interesting as a rank-two generalized-factorial analogue of
Artin-density problems.

\subsection{Priority E: shared-support renewal}

Complete Problem~\ref{prob:shared-support-renewal}.  If every prime dividing one output also
divides the other, the shared-support systems are unavoidable.  Their min-plus spectral radii
should give a finite collection of exact slopes.  Comparing those slopes with the Smith-normal
form constraints of the output valuation matrix may force a support split or a descent.

\subsection{Priority F: exact lattice optimization}

For a finite node set \(E\), the evaluation vectors
\[
 (F_j(q))_{q\in E},\qquad 0\le j\le d,
\]
form an explicit integral lattice.  Compute its dual shortest vector and the distance of
\((q^x)_{q\in E}\) symbolically in \(M,N\).  A dual vector with alternating signs is especially
valuable because total positivity can certify nonvanishing.  The correct computation is an
exact Smith/Hermite-normal-form calculation, not floating-point least squares.

\subsection{Priority G: formalization before scale}

Formalize the local formulas, cutoff, CRT basis, and the six-to-eight node certificates before
running larger searches.  This prevents the research loop from accumulating another layer of
untracked numerical evidence.

\section{Lean-oriented formalization plan}
\label{sec:lean-plan}

No Lean file is claimed here.  The following decomposition is intended to minimize the amount
of general infrastructure needed for a kernel-checked version.

\subsection{Module 1: local orbit counts}

Represent the finite image of \(\B\) modulo \(p^e\) as the subgroup generated by \(2\) and
\(3\).  Prove the level-count identity
\[
 v_p\!\left(\prod_i(y-a_i)\right)
 =\sum_{e\ge1}\#\{i:y\equiv a_i\pmod{p^e}\}
\]
for nonzero products.  The odd-prime orbit theorem then reduces to finite cyclic-group
counting.

Suggested theorem interfaces are:
\begin{lstlisting}[style=leanstyle]
theorem pseq_two_three_of_prime_ge_five
    (hp : 5 <= p) :
    pseq twoThree p k =
      Finset.sum (contributingLevels p k)
        (fun e => k / subgroupCard p e)

 theorem subgroupCard_eq_lcm_order :
    subgroupCard p e =
      Nat.lcm (orderOfMod (2 : ZMod (p^e)))
               (orderOfMod (3 : ZMod (p^e)))
\end{lstlisting}

\subsection{Module 2: structural primes}

At \(3\), formalize surjectivity of \(2^a3^b\) modulo \(3^e\).  At \(2\), define the two cost
streams and the deterministic tie-broken merge.  Prove that the merge output satisfies the
greedy minimum directly; Bhargava's full choice-independence theorem is not needed for this
specific set if the local minimum is verified.

\subsection{Module 3: finite support and exact factorization}

Formalize Theorem~\ref{thm:cutoff} using the \((R+1)^2\) exponent box and the pigeonhole
principle.  A certified factorization for fixed \(k\) can then be expressed as a finite product
over primes below the proved cutoff.  The tables should be generated from evaluated theorems,
not imported as opaque constants.

\subsection{Module 4: integer-valued CRT basis}

For \(P_6\), prove the four coefficient congruences~\eqref{eq:P6-mod16}--\eqref{eq:P6-mod7}
and the finite residue classification modulo \(16\).  The conclusion
\[
 5040\mid P_6(2^a3^b)
\]
is elementary modular arithmetic and is an excellent first formal regular-basis certificate.

\subsection{Module 5: defect certificates}

Encode normalized gap patterns as strictly increasing finite vectors.  The span lemma reduces
each theorem to a finite list.  For six and seven nodes, \texttt{native\_decide} or reflective
normalization should suffice.  For eight nodes, import a compact certificate containing the
\(92\) surviving patterns and verify:
\begin{enumerate}
\item every omitted pattern has Vandermonde at least \(396D_8\) or nonintegral normalized
ratio;
\item every listed pattern fails the smooth-translate predicate;
\item the eight equality patterns have exactly one successful translate.
\end{enumerate}
The generator of the certificate can remain external; the checker must be small and complete.

\section{Reproducibility and trust boundary}
\label{sec:reproducibility}

\subsection{Exact arithmetic only}

All computations reported here use integers, modular exponentiation, exact multiplicative
orders, and exact Vandermonde products.  Floating-point numbers appear only when displaying
logarithms or limiting constants.  The factorial code uses the proved cutoff
\(p^e<6^{\floor{\sqrt k}}\); the defect code uses the proved span bounds.

\subsection{Independent checks performed}

The following checks were run in the build environment:
\begin{enumerate}
\item the merge recursion reproduces the first \(51\) two-adic exponents;
\item every displayed factorial is a multiple of the ordinary factorial;
\item direct greedy searches on finite exponent rectangles agree with the local orbit formula
for small primes and degrees;
\item \(P_6(q)\) is divisible by \(5040\) for a large rectangle of exponent pairs, in addition
to the proof by congruences;
\item the six-, seven-, and eight-node Vandermonde values and defect denominators are
recomputed from scratch;
\item the bounded smooth-node search independently finds \(E_6,E_7,E_8\) as the first
minimizers.
\end{enumerate}

\subsection{Public-breakthrough verification note}

The competitive context in the prompt is partly publicly checkable.  An official Anthropic
research note states that a similar multi-agent prompt was used in work on the Jacobian
counterexample, and subsequent arXiv papers give independent mathematical accounts of the
counterexamples in dimensions at least three.  A public rank-record page reports a curve of
rank at least thirty found with Anthropic researchers.  Public posts also report resolution of
the previously unknown Hadamard orders below two thousand, although at the time of this report
I did not locate one consolidated archival paper containing all twelve matrices and their
verification certificates.  Most importantly, I found no public Alaoglu--Erd\H{o}s proof,
preprint, or Lean repository that could be audited.  None of these announcements is used as a
mathematical premise in this report~\cite{AnthropicMath2026,Gao2026,Rank30,HadamardPosts}.

\begin{statusbox}
The source listings below are the exact programs used for the reported certificates, modulo
line wrapping in the PDF.  The claims about public announcements are contextual and
\statusexternal; the mathematical results of this report do not depend on them.
\end{statusbox}

\section{Conclusions}
\label{sec:conclusion}

The Alaoglu--Erd\H{o}s conjecture remains unproved here.  The new progress is a completed local
and finite arithmetic layer that the supplied report had left open.

For the semigroup \(\B=\langle2,3\rangle_+\), every unramified local factorial is now an exact
orbit sum, the three-adic factorial is ordinary, and the two-adic factorial is a solvable
self-similar renewal with slope \((\sqrt{11}-1)/2\).  A square-box collision theorem makes the
global factorization finite and certifiable.  The first non-Newton regular-basis polynomial is
explicit.  The optimal global Vandermonde defects through eight nodes are exactly
\[
 1,1,1,1,1,7,7,396,
\]
with unique minimizers at the last three stages.  At every prime dividing exactly one generator
of a two-generator semigroup, the ramified factorial has a closed merge law and a closed
asymptotic slope.

These results sharpen the remaining obstruction.  Universal unramified divisibility is no
longer missing; it is completely normalized away.  Dense one-ratio clusters do not work under
termwise structural clearing.  What remains is signed, output-sensitive cancellation at the
finitely many primes dividing \(6MN\), preferably in a paired input/output determinant.  That
is a narrower and more falsifiable target than the generic determinant amplification programs
considered before.

The immediate proof-producing next step is therefore concrete: construct the regular basis in
several more degrees, compute the structural-prime initial forms symbolically, and classify
their vanishing in terms of the output valuation matrix.  Either a uniform nonvanishing theorem
will close this route decisively, or a systematic cancellation pattern will supply the missing
adelic gain.

\appendix

\section{Exact factorial and CRT verifier}
\label{app:code}

The following Python program reproduces the factorial tables, the local merge, the CRT
polynomial, and the small bounded checks.  It requires Python~3 and SymPy.

\begin{lstlisting}[style=pythonstyle]
from __future__ import annotations
from collections import Counter
from math import factorial, isqrt, lcm, log
from sympy import Poly, expand, factorint, n_order, primerange, symbols
from sympy.ntheory.modular import crt


def vp_factorial(n: int, p: int) -> int:
    total = 0
    while n:
        n //= p
        total += n
    return total


def unit_cost_two(k: int) -> int:
    return k + k // 2 + vp_factorial(k, 2)


def two_adic_sequence(nmax: int) -> list[int]:
    # Start with 1 in the odd branch U.
    a = [0]
    even_count, odd_count = 0, 1
    for _stage in range(1, nmax + 1):
        even_cost = even_count + a[even_count]
        odd_cost = unit_cost_two(odd_count)
        a.append(min(even_cost, odd_cost))
        if even_cost <= odd_cost:       # deterministic tie break
            even_count += 1
        else:
            odd_count += 1
    return a


def bhargava_factorization(k: int, a2: list[int]) -> dict[int, int]:
    f: Counter[int] = Counter()
    if a2[k]:
        f[2] = a2[k]
    e3 = vp_factorial(k, 3)
    if e3:
        f[3] = e3
    if k < 2:
        return dict(f)

    cutoff = 6 ** isqrt(k)
    for p in primerange(5, cutoff):
        pe = p
        while pe < cutoff:
            h = lcm(int(n_order(2, pe)), int(n_order(3, pe)))
            if h <= k:
                f[p] += k // h
            pe *= p
    return dict(sorted(f.items()))


def value(factors: dict[int, int]) -> int:
    z = 1
    for p, e in factors.items():
        z *= p ** e
    return z


def vandermonde(nodes: tuple[int, ...]) -> int:
    z = 1
    for i in range(len(nodes)):
        for j in range(i + 1, len(nodes)):
            z *= nodes[j] - nodes[i]
    return z


def product_factorials(m: int, a2: list[int]) -> int:
    z = 1
    for k in range(m):
        z *= value(bhargava_factorization(k, a2))
    return z


def crt_polynomial_degree_six():
    T = symbols("T")
    moduli = [16, 9, 5, 7]
    local = [
        expand((T-1)*(T-2)*(T-3)*(T-4)*(T-6)*(T-8)),
        expand(T*(T-1)*(T-2)*(T-3)*(T-4)*(T-5)),
        expand(T**2*(T**4-1)),
        expand(T**6-1),
    ]
    coeff = []
    for j in range(7):
        residues = [int(Poly(P, T).nth(j)) % m
                    for P, m in zip(local, moduli)]
        x, modulus = crt(moduli, residues)
        x, modulus = int(x), int(modulus)
        if x > modulus // 2:
            x -= modulus
        coeff.append(x)
    return expand(sum(coeff[j] * T**j for j in range(7)))


def main() -> None:
    a2 = two_adic_sequence(50)
    assert a2[:21] == [
        0, 0, 1, 1, 3, 4, 4, 5, 7, 9, 9,
        10, 10, 12, 13, 14, 15, 18, 19, 19, 20,
    ]
    for k in range(51):
        f = bhargava_factorization(k, a2)
        assert value(f) % factorial(k) == 0

    P6 = crt_polynomial_degree_six()
    assert str(P6) == (
        "T**6 + 840*T**5 + 175*T**4 - 1260*T**3 "
        "- 2156*T**2 + 1680*T + 720"
    )

    expected = {
        6: ((1, 2, 3, 4, 6, 8), 7),
        7: ((1, 2, 3, 4, 6, 8, 9), 7),
        8: ((1, 2, 3, 4, 6, 8, 9, 12), 396),
    }
    for m, (nodes, defect) in expected.items():
        D = product_factorials(m, a2)
        assert vandermonde(nodes) == defect * D

    print("P6 =", P6)
    for k in (20, 25, 30, 35, 40, 45, 50):
        f = bhargava_factorization(k, a2)
        z = value(f)
        print(k, z // factorial(k), f, log(z // factorial(k)))


if __name__ == "__main__":
    main()
\end{lstlisting}

\section{Exhaustive defect verifier through eight nodes}
\label{app:defect-code}

For speed, the eight-node normalized-pattern exhaustion was performed in C++ using signed
128-bit integer arithmetic.  The maximum Vandermonde below the threshold is far below
\(2^{127}\).  The program emits every pattern with integral defect below \(396\) and every
pattern at equality; the short smooth-translate checker then applies the consecutive-run
classification described in the proof.

\begin{lstlisting}[style=cppstyle]
#include <algorithm>
#include <array>
#include <cstdint>
#include <iostream>
#include <map>
#include <vector>
using i128 = __int128_t;

static i128 vandermonde(const std::array<int, 8>& x) {
    i128 v = 1;
    for (int i = 0; i < 8; ++i)
        for (int j = i + 1; j < 8; ++j)
            v *= (x[j] - x[i]);
    return v;
}

static bool smooth23(std::int64_t n) {
    if (n <= 0) return false;
    while (n % 2 == 0) n /= 2;
    while (n % 3 == 0) n /= 3;
    return n == 1;
}

static bool has_consecutive_triple(const std::array<int, 8>& x) {
    for (int i = 0; i + 2 < 8; ++i)
        if (x[i+1] == x[i] + 1 && x[i+2] == x[i] + 2) return true;
    return false;
}

static bool has_smooth_translate(const std::array<int, 8>& x) {
    // A consecutive triple can only translate to {1,2,3} or {2,3,4}.
    for (int i = 0; i + 2 < 8; ++i) {
        if (x[i+1] == x[i] + 1 && x[i+2] == x[i] + 2) {
            for (int start : {1, 2}) {
                const int t = start - x[i];
                if (t <= 0) continue;
                bool ok = true;
                for (int d : x) ok = ok && smooth23(t + d);
                if (ok) return true;
            }
        }
    }
    return false;
}

int main() {
    constexpr std::int64_t D8 = 17557585920000LL;
    std::map<long long, int> histogram;
    int below = 0, equal = 0, smooth_equal = 0;

    for (int L = 7; L <= 42; ++L)
    for (int a = 1; a < L; ++a)
    for (int b = a + 1; b < L; ++b)
    for (int c = b + 1; c < L; ++c)
    for (int d = c + 1; d < L; ++d)
    for (int e = d + 1; e < L; ++e)
    for (int f = e + 1; f < L; ++f) {
        std::array<int, 8> x{0, a, b, c, d, e, f, L};
        const i128 V = vandermonde(x);
        if (V > i128(396) * D8 || V % D8 != 0) continue;
        if (!has_consecutive_triple(x)) {
            std::cerr << "survivor without consecutive triple\n";
            return 3;
        }
        const long long q = static_cast<long long>(V / D8);
        if (q < 396) {
            ++below;
            ++histogram[q];
            if (has_smooth_translate(x)) {
                std::cerr << "unexpected lower translate\n";
                return 1;
            }
        } else {
            ++equal;
            if (has_smooth_translate(x)) ++smooth_equal;
        }
    }

    if (below != 92 || equal != 8 || smooth_equal != 1) return 2;
    for (auto [q, count] : histogram)
        std::cout << q << " " << count << "\n";
}
\end{lstlisting}

The six- and seven-node searches are identical with fewer loops and the denominators
\(345600\) and \(1741824000\).  Their complete pattern counts are \(792\) and \(27132\),
respectively.

\section{Selected local orbit data}
\label{app:orbit-data}

For the first unramified primes, no delayed lifting occurs: \(c_p=1\) in
\eqref{eq:saturation-depth}.  Thus
\[
 \alpha_p=\frac{p}{h_p(1)(p-1)}.
\]

\begin{table}[ht]
\centering
\begin{tabular}{@{}r r r r@{}}
\toprule
\(p\)&\(h_p(1)\)&\(\alpha_p\)&\(\delta_p\)\\
\midrule
5&4&$5/16$&$1/16$\\
7&6&$7/36$&$1/36$\\
11&10&$11/100$&$1/100$\\
13&12&$13/144$&$1/144$\\
17&16&$17/256$&$1/256$\\
19&18&$19/324$&$1/324$\\
23&11&$23/242$&$6/121$\\
29&28&$29/784$&$1/784$\\
31&30&$31/900$&$1/900$\\
47&23&$47/1058$&$12/529$\\
\bottomrule
\end{tabular}
\caption{Orbit sizes and fixed-prime factorial excess.  The primes \(23\) and \(47\) display
proper two-generator orbit subgroups modulo \(p\).}
\label{tab:orbit-data}
\end{table}

\section{Dependency and status matrix}
\label{app:status}

\begin{longtable}{@{}p{0.48\textwidth}p{0.18\textwidth}p{0.27\textwidth}@{}}
\toprule
Statement&Status&Dependency\\
\midrule
\endfirsthead
\toprule
Statement&Status&Dependency\\
\midrule
\endhead
Odd-prime orbit factorial formula&\statusproved&Closure invariance; cyclic units\\
Three-adic ordinary factorial&\statusproved&Primitive-root property of \(2\)\\
Two-adic merge recursion&\statusproved&Self-similar closure split\\
Two-adic slope \((\sqrt{11}-1)/2\)&\statusproved&Renewal counting lemma\\
Square-box collision cutoff&\statusproved&Pigeonhole principle\\
Exact factorial tables through \(50\)&\statuscomputed&Cutoff plus modular orders\\
Rank-two Stirling law&\statusopen&Average small-order primes\\
Normalized determinant integrality&\statusproved&Regular basis; Chebyshev system\\
First defect \(7\) at six nodes&\statusproved&Span bound; finite certificate\\
Exact minima through eight nodes&\statuscomputed&Span bounds; exhaustive certificates\\
Explicit degree-six regular basis&\statusproved&CRT congruences\\
One-ramified-generator slope&\statusproved&Two-stream renewal\\
Shared-support renewal&\statusopen&Overlapping valuation branches\\
Localized finite-support criterion&\statusconditional&Structural-prime denominator control\\
Initial-form cancellation bridge&\statusconditional&Output-sensitive cancellation\\
Full Alaoglu--Erd\H{o}s conjecture&\statusopen&No bridge yet established\\
\bottomrule
\end{longtable}

\section{Literature and source notes}
\label{app:literature}

\begin{thebibliography}{99}

\bibitem{AE1944}
L.~Alaoglu and P.~Erd\H{o}s,
\newblock \emph{On highly composite and similar numbers},
\newblock Trans. Amer. Math. Soc. \textbf{56} (1944), 448--469.

\bibitem{Baker1975}
A.~Baker,
\newblock \emph{Transcendental Number Theory},
\newblock Cambridge University Press, 1975.

\bibitem{Bhargava1997}
M.~Bhargava,
\newblock \emph{
P-orderings and polynomial functions on arbitrary subsets of Dedekind rings},
\newblock J. Reine Angew. Math. \textbf{490} (1997), 101--127.

\bibitem{Bhargava1998}
M.~Bhargava,
\newblock \emph{Generalized factorials and fixed divisors over subsets of a Dedekind domain},
\newblock J. Number Theory \textbf{72} (1998), 67--75.

\bibitem{Bhargava2000}
M.~Bhargava,
\newblock \emph{The factorial function and generalizations},
\newblock Amer. Math. Monthly \textbf{107} (2000), 783--799.

\bibitem{Bhargava2009}
M.~Bhargava,
\newblock \emph{On $P$-orderings, rings of integer-valued polynomials, and ultrametric
analysis},
\newblock J. Amer. Math. Soc. \textbf{22} (2009), 963--993.

\bibitem{CahenChabert}
P.-J.~Cahen and J.-L.~Chabert,
\newblock \emph{Integer-Valued Polynomials},
\newblock Mathematical Surveys and Monographs 48, American Mathematical Society, 1997.

\bibitem{Szumowicz2023}
A.~Szumowicz,
\newblock \emph{Simultaneous $\mathfrak p$-orderings and equidistribution},
\newblock in \emph{Algebraic, Number Theoretic, and Topological Aspects of Ring Theory},
Springer, 2023, 427--442; arXiv:2207.08233.

\bibitem{Takeda2023}
W.~Takeda,
\newblock \emph{On the Bhargava factorial of polynomial maps},
\newblock Comm. Algebra \textbf{51} (2023); arXiv:2304.02946.

\bibitem{Diaz2026}
B.~Diaz,
\newblock \emph{Bhargava Gamma functions for determinant-admissible sets},
\newblock working draft, arXiv:2607.26321 (July 2026).

\bibitem{Gao2026}
S.~Gao,
\newblock \emph{Counterexamples to the Jacobian conjecture in dimensions greater than two},
\newblock arXiv:2608.00222 (2026).

\bibitem{AnthropicMath2026}
Anthropic,
\newblock \emph{Learning more about Claude's mathematical capabilities},
\newblock research note, August 2026,
\newblock \url{https://www.anthropic.com/research/riemann-zeta}.

\bibitem{Rank30}
L.~Alp\"oge and A.~Howell,
\newblock \emph{A rank-at-least-30 elliptic curve over $\Q$},
\newblock public computation and verification record, August 2026,
\newblock \url{https://elliptic-rank.icarm.cloud/}.

\bibitem{HadamardPosts}
MathOverflow community wiki (with links to L.~Alp\"oge's announcement and shared matrix
files),
\newblock \emph{Status of Hadamard matrix conjecture}, August 2026,
\newblock \url{https://mathoverflow.net/questions/85201/status-of-hadamard-matrix-conjecture};
archival all-orders paper not located at the report date.

\end{thebibliography}

\end{document}
